\chapter{Многообразия и дифференциальная геометрия}
\label{chap:manifolds-differential-geometry}

Дифференциальная геометрия переносит аппарат производных и интегралов на кривые, поверхности и более общие пространства. Локально такие объекты описываются несколькими координатами, но их геометрические свойства не должны зависеть от выбранного способа параметризации. Поэтому в этой главе координатные вычисления будут отделяться от величин, характеризующих сам объект.

Первый пример --- регулярная кривая. На нём уже возникают основные идеи последующей теории: касательный вектор, метрика, длина, инвариантность относительно замены координат и кривизна как скорость изменения локального направления. Эти конструкции затем перейдут к поверхностям, многообразиям и геодезическим.

\section{Регулярные кривые}
\label{sec:regular-curves}

\subsection{Кривая, скорость и регулярность}

Пусть \(I\subseteq\R\) --- интервал. Гладкое отображение
\[
\gamma:I\to\R^3,
\qquad
 t\mapsto\gamma(t)
 =\bigl(x(t),y(t),z(t)\bigr),
\]
называется параметризованной кривой. Его образ \(\gamma(I)\) является геометрической траекторией, а число \(t\) указывает положение на ней. Один и тот же образ допускает разные параметризации: параметр может быть временем, длиной пройденного пути, углом поворота или безразмерной переменной алгоритма.

Производная
\begin{equation}
\label{eq:curve-velocity}
\dot\gamma(t)
=\dv{\gamma}{t}(t)
=\bigl(\dot x(t),\dot y(t),\dot z(t)\bigr)
\end{equation}
задаёт направление и скорость прохождения кривой. Если \(t\) измеряется в секундах, а координаты --- в метрах, то \(\dot\gamma\) имеет размерность м/с. Скалярная величина
\begin{equation}
\label{eq:curve-speed}
v(t)=\norm{\dot\gamma(t)}
\end{equation}
называется скоростью движения по траектории.

\begin{definition}
Кривая \(\gamma\in C^1(I,\R^3)\) называется регулярной, если
\[
\dot\gamma(t)\neq0
\]
для каждого \(t\in I\).
\end{definition}

Регулярность исключает остановки параметризации, в которых направление движения нельзя восстановить по первой производной. Сам образ кривой при этом может пересекать себя: условие относится к локальному поведению отображения, а не к его глобальной взаимной однозначности.

\begin{example}
Окружность радиуса \(R>0\) можно задать формулой
\[
\gamma(t)=(R\cos t,R\sin t,0).
\]
Тогда
\[
\dot\gamma(t)=(-R\sin t,R\cos t,0),
\qquad
\norm{\dot\gamma(t)}=R,
\]
поэтому кривая регулярна. Отображение \(\eta(t)=\gamma(t^3)\) имеет тот же образ вблизи точки \((R,0,0)\), но
\[
\dot\eta(0)=0.
\]
Следовательно, \(\eta\) не является регулярной параметризацией в нуле. Геометрическая окружность не приобрела особенность; выродился только выбранный параметр.
\end{example}

\subsection{Перепараметризация}

Пусть \(J\subseteq\R\) --- другой интервал, а
\[
\varphi:J\to I
\]
--- гладкое взаимно однозначное отображение с гладким обратным и \(\varphi'(u)\neq0\). Кривая
\[
\widetilde\gamma(u)=\gamma\bigl(\varphi(u)\bigr)
\]
проходит тот же геометрический путь. По правилу цепочки
\begin{equation}
\label{eq:curve-reparameterization-derivative}
\widetilde\gamma'(u)
=\dot\gamma\bigl(\varphi(u)\bigr)\varphi'(u).
\end{equation}
Если \(\varphi'(u)>0\), ориентация движения сохраняется; если \(\varphi'(u)<0\), направление прохождения меняется на противоположное. Из \eqref{eq:curve-reparameterization-derivative} видно, что численное значение скорости зависит от параметра:
\[
\norm{\widetilde\gamma'(u)}
=\norm{\dot\gamma(\varphi(u))}\abs{\varphi'(u)}.
\]
Касательная прямая, длина участка и кривизна самой траектории от такого выбора зависеть не должны. Это различие между координатной записью и геометрической величиной будет повторяться на протяжении всей главы.

\subsection{Длина дуги и натуральный параметр}

За малое приращение \(\Delta t\) вектор перемещения имеет вид
\[
\gamma(t+\Delta t)-\gamma(t)
=\dot\gamma(t)\Delta t+o(\Delta t).
\]
Его длина в первом приближении равна \(v(t)\abs{\Delta t}\). Суммирование таких малых отрезков и предельный переход приводят к следующему определению.

\begin{definition}
Длиной гладкой кривой \(\gamma:[a,b]\to\R^3\) называется число
\begin{equation}
\label{eq:curve-length}
L(\gamma)
=\int_a^b\norm{\dot\gamma(t)}\,\dd t.
\end{equation}
\end{definition}

Формула \eqref{eq:curve-length} складывает не координатные приращения, а физически пройденные расстояния. При допустимой перепараметризации \(t=\varphi(u)\) получаем
\[
\int_J
\norm{\dot\gamma(\varphi(u))}\abs{\varphi'(u)}\,\dd u
=
\int_I\norm{\dot\gamma(t)}\,\dd t.
\]
Следовательно, длина зависит от образа кривой и выбранного направления обхода не зависит.

Зафиксируем \(t_0\in I\) и определим пройденную длину
\begin{equation}
\label{eq:arc-length-function}
s(t)
=\int_{t_0}^{t}\norm{\dot\gamma(u)}\,\dd u.
\end{equation}
Для регулярной кривой
\[
\dv{s}{t}=\norm{\dot\gamma(t)}>0,
\]
поэтому \(s(t)\) строго возрастает и может использоваться как новый параметр. Если \(t=t(s)\) --- обратная функция, то кривая
\[
\widehat\gamma(s)=\gamma(t(s))
\]
имеет единичную скорость:
\begin{equation}
\label{eq:unit-speed-curve}
\norm{\dv{\widehat\gamma}{s}}
=
\norm{\dot\gamma(t(s))}\dv{t}{s}
=1.
\end{equation}
Параметр \(s\) называется натуральным параметром, или длиной дуги. Он отделяет форму траектории от закона движения по ней. В инженерной задаче геометрия задаётся функцией \(\widehat\gamma(s)\), а фактическое движение --- зависимостью \(s=s(t)\).

\subsection{Касательная и кривизна}

Для регулярной кривой единичный касательный вектор определяется формулой
\begin{equation}
\label{eq:unit-tangent}
T(t)
=\frac{\dot\gamma(t)}{\norm{\dot\gamma(t)}}.
\end{equation}
Он показывает направление траектории, не учитывая скорость её прохождения. В натуральном параметре
\[
T(s)=\dv{\gamma}{s}.
\]
Поскольку \(\inner{T(s)}{T(s)}=1\), дифференцирование даёт
\[
0
=\dv{}{s}\inner{T}{T}
=2\inner{\dv{T}{s}}{T}.
\]
Значит, изменение единичной касательной всегда перпендикулярно самой касательной: поворот направления не может происходить вдоль него.

\begin{definition}
Кривизной регулярной кривой, заданной натуральным параметром, называется величина
\begin{equation}
\label{eq:curve-curvature-definition}
\kappa(s)
=\norm{\dv{T}{s}}.
\end{equation}
В точках, где \(\kappa(s)>0\), главная нормаль определяется как
\begin{equation}
\label{eq:principal-normal}
N(s)
=\frac{1}{\kappa(s)}\dv{T}{s}.
\end{equation}

\end{definition}

Кривизна измеряет поворот касательной на единицу длины. Её размерность равна обратной длине. Для прямой \(T\) постоянно и \(\kappa=0\). Для окружности радиуса \(R\) касательная совершает полный поворот \(2\pi\) на длине \(2\pi R\), поэтому
\[
\kappa=\frac1R.
\]
Величина \(\rho=1/\kappa\) называется радиусом кривизны. Окружность радиуса \(\rho\), имеющая в рассматриваемой точке ту же касательную и ту же кривизну, является соприкасающейся окружностью: она приближает кривую до второго порядка.

Для вычислений не требуется предварительно находить натуральный параметр. Пусть
\[
v=\norm{\dot\gamma}.
\]
Тогда \(\dd/\dd s=v^{-1}\dd/\dd t\) и
\[
\dot\gamma=vT,
\qquad
\ddot\gamma=\dot v\,T+v^2\kappa N.
\]
Векторное произведение устраняет продольную компоненту \(\dot v\,T\):
\[
\dot\gamma\times\ddot\gamma
=v^3\kappa\,B,
\]
где \(B=T\times N\). Отсюда получается практическая формула
\begin{equation}
\label{eq:curve-curvature-arbitrary-parameter}
\boxed{
\kappa(t)
=
\frac{\norm{\dot\gamma(t)\times\ddot\gamma(t)}}
     {\norm{\dot\gamma(t)}^3}
}.
\end{equation}
Числитель измеряет поперечную часть ускорения, а знаменатель убирает влияние скорости параметризации.

\subsection{Кручение и репер Френе}

Если \(\kappa>0\), векторы
\[
T,
\qquad
N,
\qquad
B=T\times N
\]
образуют правый ортонормированный базис. Плоскость, натянутая на \(T\) и \(N\), называется соприкасающейся. Кривизна описывает изгиб внутри этой плоскости, но пространственная кривая может дополнительно выводить плоскость из прежнего положения.

\begin{definition}
Кручением кривой называется коэффициент \(\tau\), определяемый равенством
\begin{equation}
\label{eq:torsion-definition}
\dv{B}{s}=-\tau N.
\end{equation}
\end{definition}

Поскольку базис ортонормирован, матрица его изменения должна быть кососимметричной. Для матрицы \(Q=[\,T\;N\;B\,]\), столбцами которой служат векторы репера, выполняются формулы Френе--Серре
\begin{equation}
\label{eq:frenet-serret}
\dv{Q}{s}
=
Q
\begin{pmatrix}
0 & -\kappa & 0\\
\kappa & 0 & -\tau\\
0 & \tau & 0
\end{pmatrix}.
\end{equation}
В векторной записи это означает
\[
\dv{T}{s}=\kappa N,
\qquad
\dv{N}{s}=-\kappa T+\tau B,
\qquad
\dv{B}{s}=-\tau N.
\]
Кососимметричность гарантирует сохранение длин и взаимной ортогональности векторов репера. Кручение равно нулю для плоской кривой. Ненулевое \(\tau\) показывает скорость вращения соприкасающейся плоскости на единицу длины.

Если \(\gamma\in C^3\), а \(\dot\gamma\times\ddot\gamma\neq0\), кручение можно вычислять непосредственно по исходному параметру:
\begin{equation}
\label{eq:torsion-arbitrary-parameter}
\boxed{
\tau(t)
=
\frac{
\det\bigl(\dot\gamma(t),\ddot\gamma(t),\dddot\gamma(t)\bigr)
}{
\norm{\dot\gamma(t)\times\ddot\gamma(t)}^2
}
}.
\end{equation}
Эта формула не применяется в точках нулевой кривизны: там главная нормаль и бинормаль репера Френе не определены.

\subsection{Практическая задача: движение по винтовой линии}

Рассмотрим траекторию инструмента, который обходит цилиндр радиуса \(R\) и одновременно поднимается вдоль его оси. Пусть \(h\) --- подъём на один радиан поворота; тогда подъём за полный оборот равен \(2\pi h\). Траектория задаётся формулой
\begin{equation}
\label{eq:helix}
\gamma(\theta)
=\bigl(R\cos\theta,R\sin\theta,h\theta\bigr).
\end{equation}
Параметр \(\theta\) является углом, а не длиной и не временем. Последовательное дифференцирование даёт
\[
\begin{aligned}
\gamma'(\theta)&=(-R\sin\theta,R\cos\theta,h),\\
\gamma''(\theta)&=(-R\cos\theta,-R\sin\theta,0),\\
\gamma'''(\theta)&=(R\sin\theta,-R\cos\theta,0).
\end{aligned}
\]
Скорость относительно углового параметра постоянна, а длина участка между углами \(\theta_0\) и \(\theta_1\) равна
\[
\norm{\gamma'(\theta)}=\sqrt{R^2+h^2},
\qquad
L=\sqrt{R^2+h^2}\,\abs{\theta_1-\theta_0}.
\]
Для формул кривизны и кручения нужны величины
\[
\begin{aligned}
\gamma'\times\gamma''
&=(hR\sin\theta,-hR\cos\theta,R^2),\\
\norm{\gamma'\times\gamma''}
&=R\sqrt{R^2+h^2},\\
\det(\gamma',\gamma'',\gamma''')&=hR^2.
\end{aligned}
\]
После подстановки в \eqref{eq:curve-curvature-arbitrary-parameter} и \eqref{eq:torsion-arbitrary-parameter} получаем
\begin{equation}
\label{eq:helix-curvature-torsion}
\boxed{
\kappa=\frac{R}{R^2+h^2},
\qquad
\tau=\frac{h}{R^2+h^2}
}.
\end{equation}
При \(h=0\) винтовая линия превращается в окружность: кручение исчезает, а кривизна становится равной \(1/R\). При увеличении \(h\) траектория вытягивается вдоль оси, и её кривизна уменьшается.

Пусть теперь инструмент движется по этой траектории с физической скоростью \(u=\dd s/\dd t\). Ускорение раскладывается по касательной и главной нормали:
\begin{equation}
\label{eq:curve-acceleration-decomposition}
\frac{\dd^2\gamma}{\dd t^2}
=\dv{u}{t}T+u^2\kappa N.
\end{equation}
Первая компонента меняет модуль скорости, вторая поворачивает направление движения. Если привод допускает нормальное ускорение не более \(a_{\max}\), то
\begin{equation}
\label{eq:curvature-speed-limit}
u^2\kappa\leq a_{\max}
\qquad\Longrightarrow\qquad
u\leq\sqrt{\frac{a_{\max}}{\kappa}}.
\end{equation}
Участки большой кривизны поэтому требуют меньшей скорости даже при неизменном пределе ускорения.

Пусть \(R=0.04\) м, подъём за оборот равен \(0.02\) м, а \(a_{\max}=2\) м/с\(^2\). Тогда
\[
h\approx0.00318\ \text{м},
\quad
\kappa\approx24.84\ \text{м}^{-1},
\quad
\tau\approx1.98\ \text{м}^{-1},
\quad
u_{\max}\approx0.284\ \text{м/с}.
\]
Из равенства \(u=\sqrt{R^2+h^2}\abs{\dot\theta}\) получаем \(\abs{\dot\theta}\approx7.08\) рад/с, или \(1.13\) оборота в секунду. В реальной системе геометрический предел дополняется ограничениями на касательное ускорение, рывок и скорость приводов.

\section{Параметризованные поверхности}
\label{sec:parametric-surfaces}

Кривая использует один параметр и локально имеет одно касательное направление. Для описания поверхности нужны два независимых параметра. Они задают координатную сетку, а производная параметризации переносит малые перемещения из плоскости параметров в касательную плоскость поверхности.

\subsection{Параметризация и условие регулярности}

Пусть \(U\subseteq\R^2\) --- открытая область с координатами \((u,v)\). Гладкое отображение
\begin{equation}
\label{eq:surface-parameterization}
X:U\to\R^3,
\qquad
(u,v)\mapsto X(u,v)
=\bigl(x(u,v),y(u,v),z(u,v)\bigr),
\end{equation}
называется параметризацией поверхности. Фиксация \(v=v_0\) даёт координатную кривую \(u\mapsto X(u,v_0)\), а фиксация \(u=u_0\) --- координатную кривую \(v\mapsto X(u_0,v)\). Их скорости равны
\[
X_u=\pdv{X}{u},
\qquad
X_v=\pdv{X}{v}.
\]

Дифференциал параметризации имеет матрицу
\begin{equation}
\label{eq:surface-jacobian}
D X(u,v)
=
\begin{pmatrix}
\pdv{x}{u} & \pdv{x}{v}\\
\pdv{y}{u} & \pdv{y}{v}\\
\pdv{z}{u} & \pdv{z}{v}
\end{pmatrix}
=
\begin{bmatrix}X_u&X_v\end{bmatrix}
\in\R^{3\times2}.
\end{equation}
Для малого приращения \(h=(\Delta u,\Delta v)^{\trans}\) справедливо приближение
\begin{equation}
\label{eq:surface-linearization}
X(u+\Delta u,v+\Delta v)-X(u,v)
=D X(u,v)h+o(\norm{h}),
\end{equation}
причём
\[
D X(u,v)h=X_u\Delta u+X_v\Delta v.
\]
Именно эта линейная комбинация является первым приближением физического перемещения по поверхности.

\begin{definition}
Параметризация \(X\in C^1(U,\R^3)\) называется регулярной в точке \((u,v)\), если
\begin{equation}
\label{eq:surface-regularity}
\rankop D X(u,v)=2.
\end{equation}
Она называется регулярной на \(U\), если это условие выполнено во всех точках области.
\end{definition}

Два столбца матрицы \(D X\) должны быть линейно независимы. В \(\R^3\) это равносильно условию
\begin{equation}
\label{eq:surface-regularity-cross-product}
X_u\times X_v\neq0.
\end{equation}
Если векторное произведение обращается в нуль, два параметрических направления сливаются в одно и двумерная координатная сетка локально схлопывается.

\begin{example}
График гладкой функции \(f:\Omega\subseteq\R^2\to\R\) задаётся параметризацией
\begin{equation}
\label{eq:graph-surface}
X(u,v)=\bigl(u,v,f(u,v)\bigr).
\end{equation}
Тогда
\[
X_u=(1,0,f_u),
\qquad
X_v=(0,1,f_v),
\qquad
X_u\times X_v=(-f_u,-f_v,1).
\]
Последняя координата векторного произведения равна единице, поэтому график регулярной функции всегда является регулярной параметризованной поверхностью.
\end{example}

\subsection{Касательная плоскость и нормаль}

Пусть \(p=X(u_0,v_0)\). Любая гладкая кривая на поверхности имеет вид
\[
\gamma(t)=X\bigl(u(t),v(t)\bigr).
\]
По правилу цепочки
\begin{equation}
\label{eq:surface-curve-velocity}
\dot\gamma(t)
=X_u\dot u+X_v\dot v
=D X\begin{pmatrix}\dot u\\ \dot v\end{pmatrix}.
\end{equation}
Следовательно, все скорости кривых, проходящих через \(p\), лежат в линейной оболочке векторов \(X_u\) и \(X_v\).

\begin{definition}
Касательной плоскостью регулярной параметризованной поверхности в точке \(p=X(u_0,v_0)\) называется пространство
\begin{equation}
\label{eq:surface-tangent-plane}
T_pS
=\spanop\{X_u(u_0,v_0),X_v(u_0,v_0)\}.
\end{equation}
\end{definition}

Нормированное векторное произведение задаёт единичную нормаль
\begin{equation}
\label{eq:surface-unit-normal}
n(u,v)
=\frac{X_u\times X_v}{\norm{X_u\times X_v}}.
\end{equation}
Порядок параметров существенен: перестановка \(u\) и \(v\) меняет знак нормали. Поэтому формула \eqref{eq:surface-unit-normal} задаёт не только перпендикулярное направление, но и локальную ориентацию поверхности.

Касательная плоскость как аффинная плоскость в \(\R^3\) записывается в виде
\begin{equation}
\label{eq:surface-affine-tangent-plane}
\Pi_p
=\left\{p+aX_u+bX_v\,\middle|\,a,b\in\R\right\}.
\end{equation}
Эквивалентное уравнение через нормаль имеет вид
\[
\inner{n}{q-p}=0.
\]

Для графика \eqref{eq:graph-surface} единичная нормаль равна
\begin{equation}
\label{eq:graph-unit-normal}
n
=\frac{(-f_u,-f_v,1)}{\sqrt{1+f_u^2+f_v^2}}.
\end{equation}
Она совпадает с нормированным градиентом функции уровня
\[
F(x,y,z)=f(x,y)-z,
\qquad
F(x,y,z)=0.
\]
Так параметрическое описание связывается с неявным описанием поверхности из многомерного анализа.

\subsection{Первая фундаментальная форма}

Вектор \(h=(a,b)^{\trans}\in\R^2\) задаёт направление в плоскости параметров. После переноса на поверхность он превращается в касательный вектор
\[
D Xh=aX_u+bX_v.
\]
Его квадрат длины равен
\begin{equation}
\label{eq:induced-metric-derivation}
\begin{aligned}
\norm{D Xh}^2
&=(D Xh)^{\trans}(D Xh)\\
&=h^{\trans}(D X)^{\trans}D Xh.
\end{aligned}
\end{equation}
Отсюда возникает симметричная матрица
\begin{equation}
\label{eq:first-fundamental-form-matrix}
\boxed{
g=(D X)^{\trans}D X
}
=
\begin{pmatrix}
E&F\\
F&G
\end{pmatrix},
\end{equation}
где
\begin{equation}
\label{eq:first-fundamental-form-coefficients}
E=\inner{X_u}{X_u},
\qquad
F=\inner{X_u}{X_v},
\qquad
G=\inner{X_v}{X_v}.
\end{equation}
Матрица \(g\) называется матрицей первой фундаментальной формы, или индуцированной метрикой в координатах \((u,v)\). Она переводит координатные приращения в физические длины и углы на поверхности.

Для любого \(h\in\R^2\)
\begin{equation}
\label{eq:first-fundamental-form-quadratic}
\norm{D Xh}^2=h^{\trans}gh.
\end{equation}
Поскольку \(g=(D X)^{\trans}D X\), матрица \(g\) неотрицательно определена. При регулярной параметризации ранг \(D X\) равен двум, поэтому
\[
h^{\trans}gh>0
\quad\text{для всех }h\neq0.
\]
Иными словами, \(g\) положительно определена. Для матрицы размера \(2\times2\) это равносильно условиям
\begin{equation}
\label{eq:first-fundamental-form-positive}
E>0,
\qquad
EG-F^2>0.
\end{equation}

Коэффициент \(E\) задаёт квадрат масштаба вдоль линии \(v=\mathrm{const}\), коэффициент \(G\) --- вдоль линии \(u=\mathrm{const}\), а \(F\) измеряет их неортогональность. Координатная сетка ортогональна тогда и только тогда, когда \(F=0\).

Для графика \eqref{eq:graph-surface}
\begin{equation}
\label{eq:graph-first-fundamental-form}
g
=
\begin{pmatrix}
1+f_u^2&f_uf_v\\
f_uf_v&1+f_v^2
\end{pmatrix}.
\end{equation}
Даже если координатная область плоская и евклидова, подъём графика меняет длины и углы через производные \(f_u\) и \(f_v\).

\subsection{Длина и угол на поверхности}

Пусть кривая на поверхности задана координатами \(c(t)=(u(t),v(t))\), а её пространственный образ равен \(\gamma=X\circ c\). Из \eqref{eq:surface-curve-velocity} и \eqref{eq:first-fundamental-form-quadratic} следует
\begin{equation}
\label{eq:surface-curve-speed-metric}
\norm{\dot\gamma}^2
=
\begin{pmatrix}\dot u&\dot v\end{pmatrix}
g
\begin{pmatrix}\dot u\\\dot v\end{pmatrix}
=E\dot u^2+2F\dot u\dot v+G\dot v^2.
\end{equation}
Поэтому длина кривой вычисляется непосредственно в параметрических координатах:
\begin{equation}
\label{eq:surface-curve-length}
L(\gamma)
=\int_a^b
\sqrt{E\dot u^2+2F\dot u\dot v+G\dot v^2}\,\dd t.
\end{equation}
Коэффициенты \(E,F,G\) берутся в текущей точке \((u(t),v(t))\).

Пусть два координатных вектора \(a,b\in\R^2\) задают касательные направления \(D Xa\) и \(D Xb\). Их скалярное произведение равно
\begin{equation}
\label{eq:surface-angle-metric-product}
\inner{D Xa}{D Xb}=a^{\trans}gb.
\end{equation}
Следовательно, угол \(\alpha\) между соответствующими направлениями на поверхности определяется формулой
\begin{equation}
\label{eq:surface-angle-metric}
\cos\alpha
=
\frac{a^{\trans}gb}
{\sqrt{a^{\trans}ga}\sqrt{b^{\trans}gb}}.
\end{equation}
Так одна матрица \(g\) одновременно содержит всю локальную информацию первого порядка о длинах и углах.

\subsection{Элемент площади}

Малый прямоугольник со сторонами \(\Delta u\) и \(\Delta v\) переходит в первом приближении в параллелограмм со сторонами \(X_u\Delta u\) и \(X_v\Delta v\). Его площадь равна
\[
\norm{X_u\times X_v}\,\abs{\Delta u\Delta v}.
\]
По тождеству Грама
\begin{equation}
\label{eq:surface-area-gram}
\norm{X_u\times X_v}^2
=\norm{X_u}^2\norm{X_v}^2-\inner{X_u}{X_v}^2
=EG-F^2
=\det g.
\end{equation}
Поэтому элемент площади имеет две эквивалентные записи:
\begin{equation}
\label{eq:surface-area-element}
\boxed{
\dd A
=\norm{X_u\times X_v}\,\dd u\dd v
=\sqrt{\det g}\,\dd u\dd v
}.
\end{equation}
Площадь параметризованного патча \(S=X(U)\), покрытого без самоналожений, равна
\begin{equation}
\label{eq:parametric-surface-area}
\operatorname{Area}(S)
=\iint_U\sqrt{EG-F^2}\,\dd u\dd v.
\end{equation}

Для графика функции из \eqref{eq:graph-surface}
\[
\det g
=(1+f_u^2)(1+f_v^2)-f_u^2f_v^2
=1+f_u^2+f_v^2.
\]
Отсюда восстанавливается знакомая формула
\begin{equation}
\label{eq:graph-surface-area}
\operatorname{Area}(S)
=\iint_\Omega\sqrt{1+\norm{\grad f}^2}\,\dd u\dd v.
\end{equation}
Множитель \(\sqrt{1+\norm{\grad f}^2}\) показывает, насколько наклонный график увеличивает площадь по сравнению с его проекцией на плоскость.

\subsection{Замена параметров и локальное растяжение}

Пусть новые координаты \((\xi,\eta)\) связаны со старыми отображением
\[
(u,v)=\varphi(\xi,\eta),
\qquad
\widetilde X=X\circ\varphi.
\]
По правилу цепочки
\[
D\widetilde X=(D X)(D\varphi).
\]
Следовательно, матрица первой фундаментальной формы в новых координатах равна
\begin{equation}
\label{eq:metric-coordinate-change}
\widetilde g
=(D\widetilde X)^{\trans}D\widetilde X
=(D\varphi)^{\trans}g(D\varphi).
\end{equation}
Числа \(E,F,G\) меняются, но вычисленные по ним физические длины, углы и площади остаются прежними.

Чтобы оценивать деформацию параметрического патча относительно плоского шаблона, координаты шаблона сначала выбирают в одинаковых единицах длины. Собственные значения положительно определённой матрицы \(g\) обозначим через
\[
\lambda_1\geq\lambda_2>0.
\]
Тогда
\begin{equation}
\label{eq:surface-principal-stretches}
\sigma_1=\sqrt{\lambda_1},
\qquad
\sigma_2=\sqrt{\lambda_2}
\end{equation}
являются главными коэффициентами локального растяжения. Существуют два взаимно перпендикулярных направления в параметрической плоскости, которые растягиваются ровно в \(\sigma_1\) и \(\sigma_2\) раз.

Произведение
\begin{equation}
\label{eq:surface-area-stretch}
\sigma_1\sigma_2
=\sqrt{\lambda_1\lambda_2}
=\sqrt{\det g}
\end{equation}
задаёт коэффициент изменения площади, а отношение
\begin{equation}
\label{eq:surface-anisotropy}
q=\frac{\sigma_1}{\sigma_2}\geq1
\end{equation}
измеряет анизотропию. Значение \(q=1\) означает одинаковое растяжение во всех направлениях и сохранение малых углов. Если дополнительно \(\sigma_1=\sigma_2=1\), параметризация локально сохраняет и длины.

Эти величины применяются при контроле качества UV-развёртки, деформации сетки и параметрического проектирования. Большое \(q\) означает, что круглая малая область шаблона превращается в сильно вытянутый эллипс, а большое \(\sqrt{\det g}\) --- что площадь элемента заметно возрастает.

\subsection{Практическая задача: коническая оболочка}

Рассмотрим оболочку, радиус которой линейно меняется по высоте. Пусть \(R_0>0\) --- нижний радиус, \(H>0\) --- высота, а
\[
r(z)=R_0+kz,
\qquad
0\leq z\leq H.
\]
Чтобы обе координаты плоского шаблона измерялись в метрах, используем \(\xi=R_0\theta\), где \(\theta\) --- угол. Для полного оборота
\[
0\leq\xi\leq2\pi R_0.
\]
Параметризация оболочки имеет вид
\begin{equation}
\label{eq:conical-shell-parameterization}
X(\xi,z)
=
\left(
 r(z)\cos\frac{\xi}{R_0},
 r(z)\sin\frac{\xi}{R_0},
 z
\right).
\end{equation}
Её частные производные равны
\[
\begin{aligned}
X_\xi
&=\frac{r(z)}{R_0}
\left(-\sin\frac{\xi}{R_0},
       \cos\frac{\xi}{R_0},0\right),\\
X_z
&=\left(k\cos\frac{\xi}{R_0},
          k\sin\frac{\xi}{R_0},1\right).
\end{aligned}
\]
Их скалярное произведение равно нулю, поэтому координатная сетка ортогональна. Коэффициенты первой фундаментальной формы имеют вид
\begin{equation}
\label{eq:conical-shell-metric}
E=\left(\frac{r(z)}{R_0}\right)^2,
\qquad
F=0,
\qquad
G=1+k^2,
\end{equation}
а сама матрица равна
\[
g(\xi,z)
=
\begin{pmatrix}
\left(r(z)/R_0\right)^2&0\\
0&1+k^2
\end{pmatrix}.
\]

Векторное произведение и элемент площади вычисляются без отдельного поиска нормали:
\[
X_\xi\times X_z
=\frac{r(z)}{R_0}
\left(
\cos\frac{\xi}{R_0},
\sin\frac{\xi}{R_0},
-k
\right),
\]
\begin{equation}
\label{eq:conical-shell-area-element}
\dd A
=\frac{r(z)}{R_0}\sqrt{1+k^2}\,\dd\xi\dd z.
\end{equation}
Интегрирование по полному обороту даёт
\begin{align}
\operatorname{Area}(S)
&=\int_0^{2\pi R_0}\int_0^H
\frac{R_0+kz}{R_0}\sqrt{1+k^2}\,\dd z\dd\xi\\
&=2\pi\sqrt{1+k^2}
\left(R_0H+\frac{kH^2}{2}\right).
\label{eq:conical-shell-area}
\end{align}
Если \(R_1=R_0+kH\) --- верхний радиус, а \(\ell=H\sqrt{1+k^2}\) --- длина образующей, то формула принимает стандартный вид
\begin{equation}
\label{eq:conical-shell-area-classic}
\operatorname{Area}(S)=\pi(R_0+R_1)\ell.
\end{equation}

Поскольку матрица \(g\) диагональна, главные растяжения читаются непосредственно:
\begin{equation}
\label{eq:conical-shell-stretches}
\sigma_\xi(z)=\frac{r(z)}{R_0},
\qquad
\sigma_z=\sqrt{1+k^2}.
\end{equation}
Первое растяжение меняется по высоте и показывает расширение окружностей относительно нижнего края шаблона. Второе постоянно и возникает из-за наклона образующей. При \(k=0\) получается цилиндр, \(g=I\), и прямоугольник \((\xi,z)\) разворачивается на цилиндр без локального растяжения.

Пусть \(R_0=0.40\) м, \(R_1=0.70\) м и \(H=0.60\) м. Тогда
\[
k=\frac{R_1-R_0}{H}=0.5,
\qquad
\ell=H\sqrt{1+k^2}\approx0.671\ \text{м},
\]
\[
\operatorname{Area}(S)
=\pi(R_0+R_1)\ell
\approx2.32\ \text{м}^2.
\]
Вдоль образующей шаблон растягивается в \(\sigma_z\approx1.118\) раза. В окружном направлении растяжение возрастает от \(1\) у основания до
\[
\sigma_\xi(H)=\frac{R_1}{R_0}=1.75
\]
у верхнего края. В верхней части коэффициент изменения площади равен примерно
\[
\sigma_\xi(H)\sigma_z\approx1.96,
\]
а анизотропия ---
\[
q(H)=\frac{1.75}{1.118}\approx1.57.
\]
Такой прямоугольный способ параметризации удобен для вычислений, но сильно растягивает верхнюю часть сетки. Для изготовления оболочки из почти нерастяжимого листа нужен другой плоский шаблон --- сектор кольца, а не прямоугольник. Матрица \(g\) позволяет обнаружить эту проблему до построения физической модели.

\section{Гладкие многообразия и подмногообразия}
\label{sec:smooth-manifolds-submanifolds}

Кривая описывается одной локальной координатой, а поверхность --- двумя. Однако даже сфера не допускает одной регулярной параметризации, покрывающей её целиком без разрыва или вырождения. Поэтому вместо единой глобальной системы координат используют набор локальных систем, согласованных на пересечениях. Такой подход позволяет выполнять обычные вычисления в~\(\R^n\), не отождествляя геометрический объект с одной конкретной параметризацией.

\subsection{Локально евклидово пространство}

Пусть \(M\) --- множество точек, снабжённое понятием открытого множества. Говорят, что \(M\) локально похоже на \(\R^n\), если у каждой точки \(p\in M\) существует окрестность, которую можно взаимно однозначно и непрерывно сопоставить открытому множеству в~\(\R^n\), причём обратное отображение также непрерывно. Число \(n\) не зависит от точки и называется размерностью.

Для корректного определения многообразия добавляют два технических условия. Хаусдорфовость означает, что любые две различные точки имеют непересекающиеся окрестности; поэтому предел последовательности, если он существует, не может быть сразу двумя разными точками. Существование счётной базы означает, что все открытые множества можно строить из счётного набора базовых окрестностей. Эти требования исключают патологические пространства и не меняют привычные кривые, поверхности и множества допустимых параметров.

\begin{definition}
Топологическим многообразием размерности \(n\) называется хаусдорфово пространство со счётной базой, каждая точка которого имеет окрестность, гомеоморфную открытому подмножеству \(\R^n\).
\end{definition}

Локальное сходство с \(\R^n\) не требует вложения \(M\) в пространство большей размерности. Например, поверхность Земли можно изучать через плоские карты, не выбирая трёхмерные декартовы координаты как часть определения. Вложение полезно для визуализации, но собственная структура многообразия задаётся тем, как согласованы локальные координаты.

\subsection{Карты, атласы и переходы между координатами}

\begin{definition}
Картой на \(M\) называется пара \((U,\varphi)\), где \(U\subseteq M\) --- открытая область, а
\[
\varphi:U\to \varphi(U)\subseteq\R^n
\]
--- гомеоморфизм на открытое множество \(\varphi(U)\). Числа
\[
\varphi(p)=\bigl(x^1(p),\ldots,x^n(p)\bigr)
\]
называются локальными координатами точки \(p\).
\end{definition}

Карта переводит геометрическую задачу в обычную задачу анализа. Отображение \(\varphi\) присваивает точке координаты, а \(\varphi^{-1}\) восстанавливает точку по координатам. Параметризация поверхности из раздела~\ref{sec:parametric-surfaces} записывалась в противоположном направлении,
\[
X:\Omega\subseteq\R^2\to S.
\]
Если \(X\) взаимно однозначна на рассматриваемой области, то карта поверхности равна \(\varphi=X^{-1}\). Таким образом, параметризация и карта содержат одни и те же локальные данные, но направлены в разные стороны.

Одна карта редко покрывает весь объект. Набор карт \((U_\alpha,\varphi_\alpha)\), области которых покрывают \(M\), называется атласом. Если точка принадлежит двум областям \(U_\alpha\) и \(U_\beta\), её координаты связаны отображением
\begin{equation}
\label{eq:manifold-transition-map}
\varphi_\beta\circ\varphi_\alpha^{-1}:
\varphi_\alpha(U_\alpha\cap U_\beta)
\longrightarrow
\varphi_\beta(U_\alpha\cap U_\beta).
\end{equation}
Это отображение называется переходом между картами. Оно работает целиком внутри евклидовых пространств, поэтому его гладкость проверяется обычными производными.

Рассмотрим единичную сферу
\[
S^2=\{(x,y,z)\in\R^3:x^2+y^2+z^2=1\}.
\]
На верхней полусфере \(U_z^+=\{z>0\}\) координаты можно задать проекцией
\[
\varphi_z(x,y,z)=(x,y),
\qquad
\varphi_z^{-1}(u,v)
=\left(u,v,\sqrt{1-u^2-v^2}\right).
\]
На полусфере \(U_x^+=\{x>0\}\) возьмём
\[
\varphi_x(x,y,z)=(y,z),
\qquad
\varphi_x^{-1}(r,s)
=\left(\sqrt{1-r^2-s^2},r,s\right).
\]
На пересечении \(U_z^+\cap U_x^+\) переход имеет вид
\begin{equation}
\label{eq:sphere-chart-transition}
(\varphi_x\circ\varphi_z^{-1})(u,v)
=
\left(v,\sqrt{1-u^2-v^2}\right).
\end{equation}
Корень вычисляется только там, где его аргумент положителен, поэтому переход гладок на своей открытой области. Аналогичные карты для знаков \(x,y,z\) покрывают всю сферу. Ни одна из них не является привилегированной: геометрический объект задаётся всем согласованным атласом.

\subsection{Гладкая структура и гладкие отображения}

\begin{definition}
Две карты называются гладко согласованными, если оба перехода между ними являются бесконечно дифференцируемыми на областях определения. Атлас называется гладким, если его карты попарно гладко согласованы. Гладкая структура --- максимальный класс карт, согласованных с данным гладким атласом.
\end{definition}

Максимальность здесь не требует выписывать все возможные карты. Достаточно задать один атлас; затем к нему мысленно добавляются все карты, переходы к которым гладкие. Благодаря этому гладкость не зависит от случайного выбора минимального набора координатных областей.

Пусть \(M\) и \(N\) --- гладкие многообразия размерностей \(m\) и \(n\), а \(F:M\to N\). Выберем карту \((U,\varphi)\) около \(p\in M\) и карту \((V,\psi)\) около \(F(p)\), причём \(F(U)\subseteq V\). Координатная запись отображения равна
\begin{equation}
\label{eq:smooth-map-coordinate-representation}
F_{\psi\varphi}
=
\psi\circ F\circ\varphi^{-1}:
\varphi(U)\subseteq\R^m
\longrightarrow
\psi(V)\subseteq\R^n.
\end{equation}

\begin{definition}
Отображение \(F:M\to N\) называется гладким, если его координатная запись \eqref{eq:smooth-map-coordinate-representation} гладкая для любых допустимых карт. Гладкое взаимно однозначное отображение с гладким обратным называется диффеоморфизмом.
\end{definition}

Проверка во всех картах не создаёт дополнительной вычислительной работы. Если координатная запись гладкая в одной паре карт, то в другой паре она имеет вид
\[
\widetilde\psi\circ F\circ\widetilde\varphi^{-1}
=
(\widetilde\psi\circ\psi^{-1})
\circ
(\psi\circ F\circ\varphi^{-1})
\circ
(\varphi\circ\widetilde\varphi^{-1}).
\]
Крайние множители являются гладкими переходами между картами. По правилу цепочки гладкость средней записи сохраняется. Здесь правило цепочки из многомерного анализа обеспечивает независимость определения от координат.

Диффеоморфизм означает эквивалентность гладких геометрий: он допускает перенос функций, кривых и дифференциальных конструкций в обе стороны без потери гладкости. Например, открытый участок цилиндра после разреза вдоль образующей диффеоморфен прямоугольной области плоскости. Весь цилиндр такой одной картой не покрывается, поскольку окружная координата должна склеиваться по краям.

\subsection{Вложенные подмногообразия}

Поверхности предыдущего раздела находились внутри \(\R^3\). Более общий объект может лежать внутри произвольного многообразия \(M\) и локально выглядеть как линейная координатная плоскость.

\begin{definition}
Подмножество \(S\subseteq M\) называется вложенным подмногообразием размерности \(k\), если для каждой точки \(p\in S\) существует карта \((U,\varphi)\) многообразия \(M\), в которой
\begin{equation}
\label{eq:embedded-submanifold-normal-form}
\varphi(U\cap S)
=
\varphi(U)\cap
\bigl(\R^k\times\{0\}\bigr).
\end{equation}
Такая карта называется картой подмногообразия.
\end{definition}

В координатах \eqref{eq:embedded-submanifold-normal-form} последние \(n-k\) координат равны нулю, а первые \(k\) свободно меняются. Поэтому подмногообразие само получает карты размерности \(k\). Слово ``вложенное'' означает, что его собственная топология совпадает с топологией, унаследованной как у подмножества \(M\).

Граф гладкой функции даёт основной пример. Пусть \(f:\Omega\subseteq\R^k\to\R^m\), а
\[
\Gamma_f
=
\{(x,f(x)):x\in\Omega\}
\subseteq\R^{k+m}.
\]
Рассмотрим замену координат
\begin{equation}
\label{eq:graph-straightening-map}
\Phi(x,y)=(x,y-f(x)).
\end{equation}
Её обратное отображение равно \(\Phi^{-1}(x,z)=(x,z+f(x))\), поэтому \(\Phi\) --- диффеоморфизм. На графе \(y=f(x)\), следовательно,
\[
\Phi(x,f(x))=(x,0).
\]
Таким образом, \(\Phi\) выпрямляет граф в координатную плоскость \(\R^k\times\{0\}\), и \(\Gamma_f\) является \(k\)-мерным вложенным подмногообразием.

Этот пример показывает локальную модель общего подмногообразия. После подходящей замены координат искривлённое множество допустимых точек становится плоским, хотя исходные координаты могут давать сложные нелинейные уравнения.

\subsection{Множества уровня и регулярные значения}

Пусть \(\Omega\subseteq\R^n\) открыто, а
\[
F:\Omega\to\R^m,
\qquad
m\leq n.
\]
Для заданного \(c\in\R^m\) множество
\[
F^{-1}(c)=\{x\in\Omega:F(x)=c\}
\]
задаётся системой из \(m\) уравнений. Число уравнений само по себе ещё не гарантирует, что размерность равна \(n-m\): ограничения должны быть локально независимыми.

\begin{definition}
Значение \(c\in\R^m\) называется регулярным для \(F\), если в каждой точке \(p\in F^{-1}(c)\) линейное отображение
\[
DF(p):\R^n\to\R^m
\]
сюръективно, то есть
\begin{equation}
\label{eq:regular-value-rank-condition}
\operatorname{rank}DF(p)=m.
\end{equation}
Если \(m=1\), это условие равносильно \(\grad F(p)\neq0\).
\end{definition}

\begin{theorem}[о регулярном множестве уровня]
\label{thm:regular-level-set}
Если \(c\) --- регулярное значение гладкого отображения \(F:\Omega\subseteq\R^n\to\R^m\), то \(F^{-1}(c)\) является вложенным подмногообразием размерности
\begin{equation}
\label{eq:regular-level-set-dimension}
\dim F^{-1}(c)=n-m.
\end{equation}
\end{theorem}

Механизм теоремы совпадает с теоремой о неявной функции. Из условия ранга можно после перестановки координат разбить переменную как \((x,y)\in\R^{n-m}\times\R^m\) так, чтобы квадратная матрица \(D_yF(p)\) была обратима. Тогда уравнение
\[
F(x,y)=c
\]
локально однозначно разрешается относительно \(y\): существует гладкая функция \(y=g(x)\). Следовательно, множество уровня локально является графиком
\[
F^{-1}(c)=\{(x,g(x))\}.
\]

Можно увидеть и явное выпрямление. Введём отображение
\begin{equation}
\label{eq:level-set-straightening-map}
\Phi(x,y)=\bigl(x,F(x,y)-c\bigr).
\end{equation}
Его производная имеет блочный вид
\[
D\Phi(p)
=
\begin{pmatrix}
I_{n-m}&0\\
D_xF(p)&D_yF(p)
\end{pmatrix}.
\]
Определитель этой блочно-треугольной матрицы равен \(\det D_yF(p)\neq0\). По теореме об обратной функции \(\Phi\) является локальной заменой координат. При этом
\[
F(x,y)=c
\quad\Longleftrightarrow\quad
\Phi(x,y)=(x,0),
\]
то есть система ограничений превращается в зануление последних \(m\) координат. Формула \eqref{eq:regular-level-set-dimension} выражает простой подсчёт степеней свободы: из \(n\) переменных \(m\) независимых уравнений локально определяют \(m\) переменных через остальные.

Сфера радиуса \(R>0\) является множеством уровня функции
\[
F(x,y,z)=x^2+y^2+z^2.
\]
На уровне \(F=R^2\)
\[
\grad F(x,y,z)=2(x,y,z)\neq0,
\]
поэтому сфера --- гладкое подмногообразие размерности \(3-1=2\). Локальные параметризации из предыдущего раздела возникают именно при разрешении уравнения относительно одной из координат.

Регулярность существенна. Для
\[
F(x,y)=x^2+y^2
\]
уровень \(F^{-1}(0)=\{(0,0)\}\), но \(\grad F(0,0)=0\). Формула \(2-1=1\) здесь неверна: множество состоит из одной точки и имеет размерность ноль. В критическом уровне могут возникать самопересечения, острые точки, изменение размерности или изолированные решения, поэтому перед применением теоремы необходимо проверять ранг.

\subsection{Практическая задача: допустимые размеры оболочки}

Пусть прямоугольная оболочка имеет длину \(a>0\), ширину \(b>0\) и высоту \(h>0\). Требуется сохранить внутренний объём \(V_0\) и площадь материала \(S_0\) для коробки без верхней крышки:
\begin{equation}
\label{eq:box-design-constraints}
abh=V_0,
\qquad
ab+2h(a+b)=S_0.
\end{equation}
Параметры образуют точку \(q=(a,b,h)\in\R_{>0}^3\). Запишем ограничения как одно отображение
\begin{equation}
\label{eq:box-design-map}
F(a,b,h)
=
\begin{pmatrix}
abh-V_0\\
ab+2h(a+b)-S_0
\end{pmatrix}.
\end{equation}
Множество допустимых конструкций равно \(M=F^{-1}(0)\). Его якобиан имеет размер \(2\times3\):
\begin{equation}
\label{eq:box-design-jacobian}
DF(a,b,h)
=
\begin{pmatrix}
bh&ah&ab\\
b+2h&a+2h&2(a+b)
\end{pmatrix}.
\end{equation}

Рассмотрим проектную точку
\[
q_0=(2,1,1)\ \text{м}.
\]
Она соответствует значениям
\[
V_0=2\ \text{м}^3,
\qquad
S_0=8\ \text{м}^2.
\]
Подстановка в \eqref{eq:box-design-jacobian} даёт
\[
DF(q_0)
=
\begin{pmatrix}
1&2&2\\
3&4&6
\end{pmatrix}.
\]
Минор по переменным \((b,h)\) равен
\begin{equation}
\label{eq:box-design-regular-minor}
\det
\begin{pmatrix}
2&2\\
4&6
\end{pmatrix}
=4\neq0.
\end{equation}
Следовательно, ранг равен двум, ограничения независимы, а множество допустимых конструкций вблизи \(q_0\) является гладкой кривой размерности
\[
3-2=1.
\]
Одну величину можно менять свободно, а две другие должны подстраиваться так, чтобы оба ограничения сохранялись.

Выберем \(a\) локальным свободным параметром и будем считать \(b=b(a)\), \(h=h(a)\). Дифференцирование равенства \(F(a,b(a),h(a))=0\) даёт линейную систему для чувствительности:
\begin{equation}
\label{eq:box-design-local-sensitivity}
\begin{pmatrix}
b'(2)\\h'(2)
\end{pmatrix}
=-
\begin{pmatrix}
2&2\\4&6
\end{pmatrix}^{-1}
\begin{pmatrix}
1\\3
\end{pmatrix}
=
\begin{pmatrix}
0\\-\tfrac12
\end{pmatrix}.
\end{equation}
Следовательно, при малом изменении длины \(\Delta a\)
\[
b\approx1,
\qquad
h\approx1-\frac12\Delta a.
\]
Увеличение длины на \(0.02\) м требует в первом приближении уменьшить высоту на \(0.01\) м, а ширина изменяется лишь в более высоком порядке. Направление получено не подбором, а из локальной геометрии множества допустимых параметров.

Для общей системы ограничений порядок тот же: составить \(F(q)=0\), вычислить \(DF(q)\), проверить ранг и выбрать невырожденный минор. Он показывает, какие переменные выражаются через свободные; соответствующая линейная система даёт их чувствительность. В следующем разделе эти допустимые направления будут собраны в касательное пространство подмногообразия.

\section{Касательные и кокасательные пространства}
\label{sec:tangent-cotangent-spaces}

Касательное пространство хранит возможные мгновенные направления движения в точке независимо от координат. Двойственное ему пространство хранит линейные измерения этих направлений и приводит к дифференциалам функций.

\subsection{Касательный вектор как скорость кривой}

Пусть \(M\) --- гладкое многообразие размерности \(n\), \(p\in M\), а
\[
\gamma:(-\varepsilon,\varepsilon)\to M,
\qquad
\gamma(0)=p,
\]
--- гладкая кривая. В карте \((U,\varphi)\) около \(p\) она получает координатную запись
\[
x(t)=\varphi(\gamma(t))\in\R^n.
\]
Вектор \(\dot x(0)\) описывает скорость кривой в выбранных координатах. Совпадение координатных скоростей в одной карте по правилу цепочки сохраняется в любой другой.

\begin{definition}
Касательным вектором в точке \(p\in M\) называется класс гладких кривых через \(p\), имеющих одинаковую скорость в локальных координатах. Множество всех таких векторов обозначается через \(T_pM\) и называется касательным пространством к \(M\) в точке \(p\).
\end{definition}

Карта \(\varphi\) сопоставляет касательному вектору \(\dot x(0)\in\R^n\), поэтому \(T_pM\) является \(n\)-мерным векторным пространством. Карта задаёт его координаты, но не сам вектор.

Для вложенного многообразия \(M\subseteq\R^N\) касательный вектор является скоростью \(\dot\gamma(0)\in\R^N\). Если
\[
M=F^{-1}(c),
\qquad
F:\Omega\subseteq\R^N\to\R^m,
\]
и \(c\) --- регулярное значение, то для любой кривой \(\gamma(t)\in M\)
\[
F(\gamma(t))=c.
\]
Дифференцирование при \(t=0\) даёт
\[
DF(p)\dot\gamma(0)=0.
\]
Теорема о неявной функции даёт и обратное включение, поэтому
\begin{equation}
\label{eq:tangent-space-regular-level-set}
\boxed{
T_pM=\ker DF(p)
}.
\end{equation}
Формула \eqref{eq:tangent-space-regular-level-set} превращает геометрический вопрос о допустимых направлениях в задачу линейной алгебры.

Вернёмся к множеству допустимых размеров оболочки из предыдущего раздела. В точке \(q_0=(2,1,1)\) матрица ограничений равна
\[
DF(q_0)
=
\begin{pmatrix}
1&2&2\\
3&4&6
\end{pmatrix}.
\]
Решение системы \(DF(q_0)v=0\) имеет вид
\begin{equation}
\label{eq:box-design-tangent-space}
T_{q_0}M
=
\operatorname{span}\left\{
\begin{pmatrix}2\\0\\-1\end{pmatrix}
\right\}.
\end{equation}
Вектор \((2,0,-1)\) означает: увеличение длины на две единицы компенсируется уменьшением высоты на одну, а ширина в первом порядке не меняется. Ранее полученная чувствительность \((1,0,-1/2)\) задаёт ту же касательную прямую.

\subsection{Производная по направлению и координатный базис}

Касательный вектор можно распознать не только по кривой, но и по его действию на гладкие функции. Для \(v\in T_pM\), представленного кривой \(\gamma\), положим
\begin{equation}
\label{eq:tangent-vector-as-derivation}
v(f)
=
\left.\dv{}{t}\right|_{t=0}f(\gamma(t)),
\qquad
f\in C^\infty(M).
\end{equation}
Число \(v(f)\) является скоростью изменения \(f\) в направлении \(v\) и зависит только от первой производной кривой. По правилу произведения
\begin{equation}
\label{eq:tangent-derivation-product-rule}
v(fg)=f(p)v(g)+g(p)v(f).
\end{equation}
Линейные отображения \(C^\infty(M)\to\R\), удовлетворяющие \eqref{eq:tangent-derivation-product-rule}, называются дифференцированиями в точке \(p\) и в точности совпадают с касательными векторами.

Пусть \((x^1,\ldots,x^n)\) --- локальные координаты. Кривые, координатные скорости которых равны стандартным базисным векторам, задают координатный базис
\[
\left.\frac{\partial}{\partial x^1}\right|_p,
\ldots,
\left.\frac{\partial}{\partial x^n}\right|_p
\]
пространства \(T_pM\). Любой касательный вектор единственным образом записывается как
\begin{equation}
\label{eq:tangent-vector-coordinate-expansion}
v
=
\sum_{i=1}^n v^i
\left.\frac{\partial}{\partial x^i}\right|_p,
\qquad
v(f)
=
\sum_{i=1}^n v^i\pdv{(f\circ\varphi^{-1})}{x^i}\bigl(\varphi(p)\bigr).
\end{equation}
При замене координат \(y=y(x)\) базисные векторы связаны правилом цепочки:
\begin{equation}
\label{eq:coordinate-basis-transformation}
\left.\frac{\partial}{\partial x^i}\right|_p
=
\sum_{j=1}^n
\pdv{y^j}{x^i}(p)
\left.\frac{\partial}{\partial y^j}\right|_p.
\end{equation}
Компоненты \(v^i\) меняются согласованно с базисом, поэтому действие \(v(f)\) остаётся неизменным.

\subsection{Кокасательное пространство и дифференциал функции}

Двойственное пространство из линейной алгебры даёт линейные измерители касательных направлений.

\begin{definition}
Кокасательным пространством к \(M\) в точке \(p\) называется двойственное пространство
\begin{equation}
\label{eq:cotangent-space-definition}
T_p^*M=(T_pM)^*.
\end{equation}
Его элементы называются ковекторами.
\end{definition}

Координатному базису касательных векторов соответствует двойственный базис
\[
\dd x^1\big|_p,\ldots,\dd x^n\big|_p,
\]
определяемый равенствами
\begin{equation}
\label{eq:cotangent-dual-basis}
\dd x^i\big|_p
\left(
\left.\frac{\partial}{\partial x^j}\right|_p
\right)
=\delta^i_j.
\end{equation}
Если \(\alpha\in T_p^*M\), то
\[
\alpha
=
\sum_{i=1}^n\alpha_i\,\dd x^i\big|_p,
\qquad
\alpha(v)=\sum_{i=1}^n\alpha_i v^i.
\]
Вектор задаёт движение, а ковектор возвращает измеряемое вдоль него число.

Для гладкой функции \(f:M\to\R\) её дифференциал в точке \(p\) определяется формулой
\begin{equation}
\label{eq:function-differential-on-manifold}
\dd f_p(v)=v(f).
\end{equation}
Следовательно, \(\dd f_p\in T_p^*M\), а в координатах
\begin{equation}
\label{eq:function-differential-coordinates}
\dd f_p
=
\sum_{i=1}^n
\pdv{(f\circ\varphi^{-1})}{x^i}\bigl(\varphi(p)\bigr)
\dd x^i\big|_p.
\end{equation}
Дифференциал существует без скалярного произведения. После выбора метрики ему будет соответствовать градиент по правилу
\[
\dd f_p(v)=g_p(\operatorname{grad}f(p),v).
\]
До введения метрики дифференциал и градиент отождествлять нельзя.

Для системы ограничений \(F=(F^1,\ldots,F^m)\) касательное пространство регулярного уровня можно записать через ковекторы:
\begin{equation}
\label{eq:tangent-space-annihilated-by-constraints}
T_pF^{-1}(c)
=
\ker\dd F^1_p\cap\cdots\cap\ker\dd F^m_p.
\end{equation}
В примере оболочки
\[
\dd F^1_{q_0}=\dd a+2\dd b+2\dd h,
\qquad
\dd F^2_{q_0}=3\dd a+4\dd b+6\dd h.
\]
Оба ковектора обращаются в нуль на \((2,0,-1)\): касательное движение не нарушает ограничения в первом порядке. Их линейная оболочка называется конормальным пространством.

\subsection{Касательное отображение, pushforward и pullback}

Пусть \(F:M\to N\) --- гладкое отображение и \(p\in M\). Если \(v\in T_pM\) представлен кривой \(\gamma\), то композиция \(F\circ\gamma\) является кривой в \(N\), проходящей через \(F(p)\). Её скорость определяет линейное отображение
\begin{equation}
\label{eq:tangent-map-definition}
DF(p):T_pM\to T_{F(p)}N,
\qquad
DF(p)v
=
\left.\dv{}{t}\right|_{t=0}F(\gamma(t)).
\end{equation}
Это дифференциал, или касательное отображение; запись \(F_{*p}v=DF(p)v\) называют pushforward.

В картах \((U,\varphi)\) на \(M\) и \((V,\psi)\) на \(N\) матрицей \(DF(p)\) служит обычный якобиан
\begin{equation}
\label{eq:tangent-map-coordinate-jacobian}
D\bigl(\psi\circ F\circ\varphi^{-1}\bigr)\bigl(\varphi(p)\bigr).
\end{equation}
Это знакомый дифференциал из многомерного анализа, записанный независимо от координат. Для композиции выполняется
\begin{equation}
\label{eq:tangent-map-chain-rule}
D(G\circ F)(p)
=
DG(F(p))\circ DF(p).
\end{equation}

Параметризация поверхности \(X:\Omega\subseteq\R^2\to S\) даёт непосредственный пример pushforward:
\begin{equation}
\label{eq:parametrization-pushforward-basis}
DX(u,v)e_1=X_u(u,v),
\qquad
DX(u,v)e_2=X_v(u,v).
\end{equation}
Якобиан параметризации переносит координатные направления в касательную плоскость, а \(g=(DX)^{\trans}DX\) измеряет их скалярные произведения.

Перенос функций и ковекторов идёт в обратную сторону. Для функции \(h:N\to\R\) определяют pullback
\begin{equation}
\label{eq:function-pullback}
F^*h=h\circ F.
\end{equation}
Если \(\alpha\in T_{F(p)}^*N\), то её pullback является ковектором в \(p\):
\begin{equation}
\label{eq:covector-pullback}
(F^*\alpha)_p(v)
=
\alpha_{F(p)}\bigl(DF(p)v\bigr).
\end{equation}
Если матрица \(DF(p)\) равна \(A\), то \(v\mapsto Av\), а коэффициенты ковектора переносятся как \(\alpha\mapsto A^{\trans}\alpha\). Этим объясняется транспонированный якобиан в обратном распространении производных.

Для скалярной функции \(h\) формулы согласованы:
\begin{equation}
\label{eq:differential-pullback-chain-rule}
\dd(F^*h)_p
=
F^*(\dd h)_{p},
\qquad
\dd(h\circ F)_p(v)
=
\dd h_{F(p)}\bigl(DF(p)v\bigr).
\end{equation}
Это равенство позволит переносить формы с поверхности на параметрическую область.

\subsection{Векторные поля и потоки}

Для описания движения допустимое направление выбирают гладко в каждой точке.

\begin{definition}
Векторным полем на \(M\) называется гладкое правило
\[
p\longmapsto X_p\in T_pM.
\]
В локальных координатах оно имеет вид
\begin{equation}
\label{eq:vector-field-local-form}
X
=
\sum_{i=1}^n X^i(x)
\frac{\partial}{\partial x^i},
\end{equation}
где коэффициенты \(X^i\) являются гладкими функциями.
\end{definition}

Векторное поле действует на функцию как дифференциальный оператор первого порядка:
\[
X(f)(p)=X_p(f).
\]
Кривая \(\gamma(t)\) называется интегральной кривой поля \(X\), если
\begin{equation}
\label{eq:integral-curve-vector-field}
\dot\gamma(t)=X_{\gamma(t)}.
\end{equation}
В координатах \eqref{eq:integral-curve-vector-field} превращается в систему обыкновенных дифференциальных уравнений
\begin{equation}
\label{eq:vector-field-coordinate-ode}
\dot x^i(t)=X^i(x(t)),
\qquad
i=1,\ldots,n.
\end{equation}
Так поле задаёт автономную динамическую систему на многообразии.

Решение с начальным условием \(\gamma(0)=p\) обозначим через
\[
\Phi_t(p)=\gamma_p(t).
\]
Отображение \(\Phi_t\) называется потоком. Единственность решения даёт
\begin{equation}
\label{eq:flow-group-property}
\Phi_0=\operatorname{id}_M,
\qquad
\Phi_{t+s}=\Phi_t\circ\Phi_s
\end{equation}
для всех времён, для которых обе стороны определены. При существовании решений для всех \(t\in\R\) поле называется полным.

Рассмотрим окружность
\[
S^1=\{(\cos\theta,\sin\theta):\theta\in\R\}.
\]
Единичное касательное направление равно
\[
e_\theta=(-\sin\theta,\cos\theta).
\]
Поле постоянного вращения с угловой скоростью \(\omega\) задаётся формулой
\[
X_{(\cos\theta,\sin\theta)}=\omega e_\theta.
\]
Уравнение интегральной кривой сводится к \(\dot\theta=\omega\), поэтому
\begin{equation}
\label{eq:circle-rotation-flow}
\Phi_t(\cos\theta,\sin\theta)
=
\bigl(\cos(\theta+\omega t),\sin(\theta+\omega t)\bigr).
\end{equation}
Поле задаёт мгновенную скорость, а поток --- конечное положение; это прямой мост к ОДУ и динамическим системам.

\subsection{Скобка Ли как несовместимость двух движений}

Разность двух порядков применения полей к функциям снова является векторным полем.

\begin{definition}
Скобкой Ли векторных полей называется поле
\begin{equation}
\label{eq:lie-bracket-definition}
[X,Y](f)=X(Y(f))-Y(X(f)).
\end{equation}
В координатах, если \(X=\sum_iX^i\partial_i\) и \(Y=\sum_iY^i\partial_i\), то
\begin{equation}
\label{eq:lie-bracket-coordinates}
[X,Y]
=
\sum_{i=1}^n
\left(
\sum_{j=1}^n
X^j\pdv{Y^i}{x^j}
-
Y^j\pdv{X^i}{x^j}
\right)
\frac{\partial}{\partial x^i}.
\end{equation}
\end{definition}

Скобка измеряет, насколько два малых движения не коммутируют. Рассмотрим на \(\R^2\) поля
\[
X=\frac{\partial}{\partial x},
\qquad
Y=x\frac{\partial}{\partial y}.
\]
Первое поле сдвигает точку по горизонтали, а второе --- по вертикали со скоростью, равной текущей координате \(x\). Их потоки равны
\[
\Phi_t^X(x,y)=(x+t,y),
\qquad
\Phi_s^Y(x,y)=(x,y+sx).
\]
В двух порядках получаем
\begin{align}
\Phi_t^X\circ\Phi_s^Y(x,y)
&=(x+t,y+sx),\\
\Phi_s^Y\circ\Phi_t^X(x,y)
&=(x+t,y+s(x+t)).
\end{align}
Второй результат выше первого на вектор \((0,st)\). Вычисление по \eqref{eq:lie-bracket-coordinates} даёт
\begin{equation}
\label{eq:lie-bracket-noncommuting-example}
[X,Y]=\frac{\partial}{\partial y}.
\end{equation}
Скобка указывает дополнительное направление, возникающее из-за порядка управляющих воздействий; в общем случае \([X,Y]=0\) равносильно локальной перестановочности потоков.

\section{Дифференциальные формы на многообразиях}
\label{sec:differential-forms-manifolds}

В главе о многомерном анализе дифференциальные формы уже связали градиент, ротор, дивергенцию и классические интегральные формулы в пространстве \(\R^3\). Теперь тот же аппарат переносится на произвольное гладкое многообразие. Ковектор измеряет одно касательное направление; форма степени \(k\) одновременно измеряет ориентированный \(k\)-мерный элемент.

\subsection{Форма как многолинейное измерение}

Пусть \(M\) --- гладкое многообразие размерности \(n\), а \(p\in M\). Обозначим через
\[
\Lambda^k(T_p^*M)
\]
пространство кососимметрических \(k\)-линейных отображений
\[
\omega_p:\underbrace{T_pM\times\cdots\times T_pM}_{k\ \text{аргументов}}\to\R.
\]
Кососимметричность означает, что перестановка двух аргументов меняет знак. Поэтому \(\omega_p(v_1,\ldots,v_k)=0\), если аргументы линейно зависимы: вырожденный параллелепипед не имеет ориентированного \(k\)-мерного объёма.

\begin{definition}
Гладкой дифференциальной \(k\)-формой на \(M\) называется гладко зависящий от точки выбор
\[
p\longmapsto \omega_p\in\Lambda^k(T_p^*M).
\]
Пространство таких форм обозначается через \(\Omega^k(M)\).
\end{definition}

При \(k=0\) форма является гладкой функцией, то есть \(\Omega^0(M)=C^\infty(M)\). При \(k=1\) получаем поле ковекторов из предыдущего раздела. Формы степени \(k>n\) тождественно равны нулю, поскольку в \(n\)-мерном касательном пространстве любые \(k\) векторов линейно зависимы.

В локальных координатах \((x^1,\ldots,x^n)\) форма записывается как
\begin{equation}
\label{eq:k-form-local-coordinates}
\omega
=
\sum_{1\le i_1<\cdots<i_k\le n}
\omega_{i_1\ldots i_k}(x)\,
\dd x^{i_1}\wedge\cdots\wedge\dd x^{i_k}.
\end{equation}
Коэффициенты \(\omega_{i_1\ldots i_k}\) зависят от координат, но значение \(\omega_p(v_1,\ldots,v_k)\) является геометрическим числом. Число независимых коэффициентов равно
\[
\dim\Lambda^k(T_p^*M)=\binom{n}{k}.
\]
На поверхности, где \(n=2\), существуют только функции, \(1\)-формы и \(2\)-формы. Любая \(2\)-форма локально имеет вид \(a(u,v)\,\dd u\wedge\dd v\).

\subsection{Внешнее произведение}

Внешнее произведение соединяет измерения разных размерностей. Для \(1\)-форм \(\alpha\) и \(\beta\)
\begin{equation}
\label{eq:wedge-two-one-forms-manifold}
(\alpha\wedge\beta)(u,v)
=
\alpha(u)\beta(v)-\alpha(v)\beta(u).
\end{equation}
Правая часть является ориентированной площадью образа параллелограмма после линейного отображения
\[
w\longmapsto \bigl(\alpha(w),\beta(w)\bigr).
\]
Для форм произвольных степеней внешнее произведение
\[
\wedge:\Omega^k(M)\times\Omega^\ell(M)\to\Omega^{k+\ell}(M)
\]
задаётся многолинейностью и кососимметризацией. Оно удовлетворяет правилу
\begin{equation}
\label{eq:wedge-graded-commutativity}
\alpha\wedge\beta
=(-1)^{k\ell}\beta\wedge\alpha,
\qquad
\alpha\in\Omega^k(M),\ \beta\in\Omega^\ell(M).
\end{equation}
В частности, \(\alpha\wedge\alpha=0\) для любой \(1\)-формы \(\alpha\). В координатах вычисление сводится к раскрытию скобок с учётом
\[
\dd x^i\wedge\dd x^j=-\dd x^j\wedge\dd x^i,
\qquad
\dd x^i\wedge\dd x^i=0.
\]

\begin{example}
Пусть на двумерном многообразии
\[
\alpha=a\,\dd u+b\,\dd v,
\qquad
\beta=c\,\dd u+d\,\dd v.
\]
Тогда
\begin{equation}
\label{eq:wedge-two-one-forms-determinant}
\alpha\wedge\beta
=(ad-bc)\,\dd u\wedge\dd v.
\end{equation}
Коэффициент \(ad-bc\) является определителем матрицы, составленной из коэффициентов форм. Внешнее произведение тем самым продолжает геометрический смысл определителя как ориентированного коэффициента площади.
\end{example}

\subsection{Pullback формы}

Пусть \(F:M\to N\) --- гладкое отображение, а \(\omega\in\Omega^k(N)\). Касательное отображение переносит векторы вперёд,
\[
DF(p):T_pM\to T_{F(p)}N,
\]
поэтому форма естественно переносится назад.

\begin{definition}
Pullback формы \(\omega\) отображением \(F\) определяется равенством
\begin{equation}
\label{eq:k-form-pullback-definition}
(F^*\omega)_p(v_1,\ldots,v_k)
=
\omega_{F(p)}\bigl(DF(p)v_1,\ldots,DF(p)v_k\bigr).
\end{equation}
\end{definition}

При \(k=0\) это обычная композиция \(F^*f=f\circ F\), а при \(k=1\) формула совпадает с pullback ковектора из предыдущего раздела. Операция сохраняет внешнее произведение:
\begin{equation}
\label{eq:pullback-preserves-wedge}
F^*(\alpha\wedge\beta)
=F^*\alpha\wedge F^*\beta.
\end{equation}

Координатный рецепт следует непосредственно из правила цепочки. Если \(y^a\) --- координаты на \(N\), а
\[
y^a=F^a(x^1,\ldots,x^n),
\]
то
\begin{equation}
\label{eq:pullback-coordinate-differentials}
F^*(\dd y^a)
=\dd(F^a)
=
\sum_{i=1}^n\pdv{F^a}{x^i}\,\dd x^i.
\end{equation}
Далее в форму подставляются выражения \eqref{eq:pullback-coordinate-differentials}, раскрываются внешние произведения и собираются коэффициенты.

Для параметризации поверхности
\[
X:U\subseteq\R^2\to S
\]
pullback переводит форму на искривлённой поверхности в обычную форму на плоской области \(U\). Именно поэтому интеграл по поверхности вычисляется двойным интегралом по параметрам.

\subsection{Внешний дифференциал}

Внешний дифференциал продолжает обычный дифференциал функции и повышает степень формы на единицу:
\[
\dd:\Omega^k(M)\to\Omega^{k+1}(M).
\]
В координатах для формы \eqref{eq:k-form-local-coordinates}
\begin{equation}
\label{eq:exterior-derivative-manifold-coordinates}
\dd\omega
=
\sum_{i_1<\cdots<i_k}
\sum_{j=1}^n
\pdv{\omega_{i_1\ldots i_k}}{x^j}
\dd x^j\wedge\dd x^{i_1}\wedge\cdots\wedge\dd x^{i_k}.
\end{equation}
Хотя правая часть записана в координатах, результат при их замене остаётся той же формой. Это следует из правила цепочки и сокращения слагаемых со вторыми производными благодаря кососимметрии.

Внешний дифференциал удовлетворяет двум основным тождествам:
\begin{align}
\label{eq:exterior-derivative-leibniz-manifold}
\dd(\alpha\wedge\beta)
&=
\dd\alpha\wedge\beta
+(-1)^k\alpha\wedge\dd\beta,
\qquad \alpha\in\Omega^k(M),\\
\label{eq:exterior-derivative-square-zero-manifold}
\dd(\dd\alpha)&=0.
\end{align}
Кроме того, внешний дифференциал согласован с pullback:
\begin{equation}
\label{eq:pullback-commutes-exterior-derivative}
\dd(F^*\omega)=F^*(\dd\omega).
\end{equation}
Формула \eqref{eq:pullback-commutes-exterior-derivative} позволяет вычислять производную формы после перехода к удобным координатам. Равенство \(\dd^2=0\) в \(\R^3\) содержит знакомые тождества \(\operatorname{rot}(\grad f)=0\) и \(\operatorname{div}(\operatorname{rot}F)=0\).

\subsection{Ориентация и формы старшей степени}

Интеграл формы должен различать два направления обхода. Ориентация \(n\)-мерного многообразия согласованно выбирает, какие упорядоченные базисы каждого \(T_pM\) считать положительными. В ориентированном атласе якобианы переходов между картами имеют положительные определители.

Ненулевая \(n\)-форма \(\mu\) задаёт ориентацию условием
\[
\mu_p(v_1,\ldots,v_n)>0
\]
на положительных базисах. Такая форма называется формой ориентации. В положительной карте она имеет вид
\begin{equation}
\label{eq:orientation-form-local}
\mu=\rho(x)\,\dd x^1\wedge\cdots\wedge\dd x^n,
\qquad
\rho(x)>0.
\end{equation}
Смена ориентации меняет знак \(\mu\). Ориентация задаёт знак, но ещё не задаёт канонический геометрический объём: плотность \(\rho\) выбирается дополнительно. В следующем разделе риманова метрика определит естественную форму объёма.

Если \(M\) имеет край, ориентация \(\partial M\) согласуется с ориентацией \(M\) правилом «сначала внешняя нормаль». Базис \((v_1,\ldots,v_{n-1})\) пространства \(T_p\partial M\) считается положительным, когда
\[
(\nu_{\mathrm{out}},v_1,\ldots,v_{n-1})
\]
является положительным базисом \(T_pM\). Для плоской области со стандартной ориентацией это даёт обход границы против часовой стрелки.

\subsection{Интегрирование формы}

Пусть \(M\) --- ориентированное \(k\)-мерное многообразие, параметризованное одной положительно ориентированной картой
\[
X:U\subseteq\R^k\to M.
\]
Для компактно поддержанной формы \(\omega\in\Omega^k(M)\) определяют
\begin{equation}
\label{eq:form-integral-by-pullback}
\int_M\omega
:=
\int_U X^*\omega.
\end{equation}
Если параметризация заменена другой положительно ориентированной параметризацией, формула замены переменных оставляет значение интеграла неизменным. При смене ориентации знак интеграла меняется. Для многообразия, не покрываемого одной картой, его разбивают на конечное число ориентированных кусков и складывают интегралы; согласованность на перекрытиях обеспечивается той же формулой замены переменных.

Интегрировать можно только форму той же степени, что и размерность области. \(1\)-форма интегрируется вдоль кривой, \(2\)-форма --- по поверхности, \(n\)-форма --- по \(n\)-мерной области. Pullback в \eqref{eq:form-integral-by-pullback} автоматически подставляет касательные направления параметризации и учитывает ориентированный коэффициент площади или объёма.

\begin{theorem}[Общая формула Стокса]
\label{thm:general-stokes-manifold}
Пусть \(M\) --- компактное ориентированное \(k\)-мерное гладкое многообразие с краем, а \(\omega\in\Omega^{k-1}(M)\). Тогда
\begin{equation}
\label{eq:general-stokes-manifold}
\boxed{
\int_M\dd\omega
=
\int_{\partial M}\omega
}.
\end{equation}
Ориентация края согласована с ориентацией \(M\) правилом внешней нормали.
\end{theorem}

Формула \eqref{eq:general-stokes-manifold} утверждает, что суммарное внутреннее изменение формы определяется её значениями на границе. При \(k=1\) она даёт формулу Ньютона--Лейбница, при \(k=2\) --- формулу Грина и классическую формулу Стокса, при \(k=3\) --- формулу Гаусса--Остроградского. Эти результаты отличаются размерностью, но используют один и тот же оператор \(\dd\).

\subsection{Поток через полусферу: pullback и Стокс}

Рассмотрим поле скоростей
\begin{equation}
\label{eq:radial-expansion-field}
u(x,y,z)=\alpha(x,y,z),
\qquad
\alpha>0,
\end{equation}
где координаты измеряются в метрах, а \(\alpha\) --- в \(\text{с}^{-1}\). Поле описывает равномерное радиальное расширение. Его форме потока соответствует \(2\)-форма
\begin{equation}
\label{eq:radial-field-flux-form}
\eta
=
\alpha x\,\dd y\wedge\dd z
+
\alpha y\,\dd z\wedge\dd x
+
\alpha z\,\dd x\wedge\dd y.
\end{equation}
Найдём поток через верхнюю полусферу радиуса \(R\), ориентированную наружу. Параметризация
\[
X(\theta,\varphi)
=
\bigl(
R\sin\theta\cos\varphi,
R\sin\theta\sin\varphi,
R\cos\theta
\bigr)
\]
задана при
\[
0\le\theta\le\frac{\pi}{2},
\qquad
0\le\varphi\le2\pi.
\]
Её касательные векторы равны
\begin{align*}
X_\theta
&=
R(\cos\theta\cos\varphi,
\cos\theta\sin\varphi,-\sin\theta),\\
X_\varphi
&=
R(-\sin\theta\sin\varphi,
\sin\theta\cos\varphi,0).
\end{align*}
Их векторное произведение
\begin{equation}
\label{eq:hemisphere-oriented-area-vector}
X_\theta\times X_\varphi
=
R^2\sin\theta
(\sin\theta\cos\varphi,
\sin\theta\sin\varphi,
\cos\theta)
\end{equation}
направлено наружу. Подстановка в определение pullback даёт
\begin{align}
X^*\eta
&=
\inner{u(X)}{X_\theta\times X_\varphi}
\,\dd\theta\wedge\dd\varphi\notag\\
&=
\alpha R^3\sin\theta
\,\dd\theta\wedge\dd\varphi.
\label{eq:hemisphere-flux-pullback}
\end{align}
Следовательно,
\begin{equation}
\label{eq:hemisphere-flux-direct}
\int_{S_+}\eta
=
\int_0^{2\pi}\int_0^{\pi/2}
\alpha R^3\sin\theta\,\dd\theta\,\dd\varphi
=
2\pi\alpha R^3.
\end{equation}
Результат имеет размерность \(\text{м}^3/\text{с}\), как и объёмный поток.

Проверим вычисление общей формулой Стокса. Из \eqref{eq:radial-field-flux-form}
\begin{equation}
\label{eq:radial-flux-form-exterior-derivative}
\dd\eta
=3\alpha\,\dd x\wedge\dd y\wedge\dd z.
\end{equation}
Пусть \(B_+\) --- верхний полушар. Его граница состоит из полусферы \(S_+\) и диска \(D\) в плоскости \(z=0\). На диске нормальная компонента поля равна нулю, поэтому \(\int_D\eta=0\). По теореме Стокса
\begin{align}
\int_{S_+}\eta
&=
\int_{\partial B_+}\eta
=
\int_{B_+}\dd\eta\notag\\
&=
3\alpha\operatorname{Vol}(B_+)
=
3\alpha\frac{2\pi R^3}{3}
=
2\pi\alpha R^3.
\label{eq:hemisphere-flux-stokes-check}
\end{align}
Прямой расчёт использует локальные касательные векторы параметризации, а формула Стокса заменяет поверхностный интеграл объёмным. Совпадение показывает, что pullback, ориентация и внешний дифференциал являются не отдельными формальностями, а согласованным вычислительным аппаратом.

\section{Риманова метрика и оптимизация на многообразиях}
\label{sec:riemannian-metric-optimization}

Касательное пространство задаёт допустимые направления движения, но само по себе не определяет их длины и углы между ними. Для этого на каждом касательном пространстве выбирают скалярное произведение. Если выбор гладко меняется от точки к точке, многообразие получает внутреннюю геометрию: можно измерять длины кривых, площади и объёмы, а также определять направление наискорейшего изменения функции, не выходя из множества ограничений.

\subsection{Риманова метрика}

Пусть $M$ --- гладкое многообразие размерности $n$. Для каждой точки $p\in M$ касательное пространство $T_pM$ является $n$-мерным векторным пространством.

\begin{definition}
Римановой метрикой на $M$ называется семейство скалярных произведений
\[
g_p:T_pM\times T_pM\to\R,
\qquad p\in M,
\]
которое гладко зависит от точки $p$.
\end{definition}

Для фиксированной точки $p$ отображение $g_p$ линейно по каждому аргументу, симметрично и положительно определено:
\begin{equation}
\label{eq:riemannian-metric-properties}
g_p(v,w)=g_p(w,v),
\qquad
g_p(v,v)>0\quad\text{при }v\neq0.
\end{equation}
Отсюда определяются норма и угол:
\begin{equation}
\label{eq:riemannian-norm-angle}
\norm{v}_{g,p}=\sqrt{g_p(v,v)},
\qquad
\cos\angle_g(v,w)
=
\frac{g_p(v,w)}{\norm{v}_{g,p}\norm{w}_{g,p}}.
\end{equation}
В евклидовом пространстве $M=\R^n$ стандартная метрика не зависит от точки и совпадает с обычным скалярным произведением:
\[
g_p(v,w)=v^{\trans}w.
\]
На общем многообразии пространства $T_pM$ и $T_qM$ различны, поэтому метрика сравнивает векторы только в одной точке. Сравнение векторов из разных касательных пространств потребует параллельного переноса, который будет введён в следующем разделе.

\subsection{Координатная матрица метрики}

Пусть $(x^1,\ldots,x^n)$ --- локальные координаты и
\[
\partial_i:=\pdv{}{x^i}
\]
--- соответствующий базис $T_pM$. Коэффициенты метрики определяются формулами
\begin{equation}
\label{eq:metric-coefficients}
g_{ij}(p)=g_p(\partial_i,\partial_j).
\end{equation}
Они образуют симметричную положительно определённую матрицу
\[
g(p)=\bigl(g_{ij}(p)\bigr)_{i,j=1}^n\in\R^{n\times n}.
\]
Если
\[
v=\sum_{i=1}^n v^i\partial_i,
\qquad
w=\sum_{j=1}^n w^j\partial_j,
\]
то
\begin{equation}
\label{eq:metric-coordinate-quadratic-form}
g_p(v,w)
=
\sum_{i,j=1}^n g_{ij}(p)v^iw^j
=
v^{\trans}g(p)w.
\end{equation}
Таким образом, риманова метрика является гладко меняющейся положительно определённой квадратичной формой. Связь с линейной алгеброй непосредственна: собственные значения $g(p)$ показывают, насколько единичные координатные направления растягиваются или сжимаются при измерении геометрической длины.

При замене координат $x=x(y)$ базис меняется по правилу цепочки. Если
\[
J=\pdv{x}{y}
\]
--- матрица перехода от $y$-координат к $x$-координатам, то новая матрица метрики равна
\begin{equation}
\label{eq:riemannian-metric-coordinate-change}
\widetilde g(y)=J(y)^{\trans}g(x(y))J(y).
\end{equation}
Матрица зависит от координат, но число $g_p(v,w)$ не меняется: одновременно преобразуются и матрица, и координатные столбцы векторов.

\subsection{Индуцированная метрика, длина и объём}

Пусть $M\subseteq\R^m$ --- вложенное многообразие. Евклидово скалярное произведение можно ограничить на каждое касательное пространство:
\begin{equation}
\label{eq:induced-metric-submanifold}
g_p(v,w)=v^{\trans}w,
\qquad
v,w\in T_pM\subseteq\R^m.
\end{equation}
Полученная метрика называется индуцированной. В локальной параметризации
\[
X:U\subseteq\R^n\to M\subseteq\R^m
\]
координатные касательные векторы являются столбцами матрицы $DX$. Поэтому
\begin{equation}
\label{eq:induced-metric-jacobian-v06}
\boxed{g=(DX)^{\trans}DX}.
\end{equation}
Для поверхности $X(u,v)$ эта формула совпадает с матрицей первой фундаментальной формы из раздела~\ref{sec:parametric-surfaces}:
\[
g=
\begin{pmatrix}
E&F\\
F&G
\end{pmatrix}.
\]
Здесь $E,F,G$ --- традиционные коэффициенты первой фундаментальной формы, а строчная буква $g$ обозначает всю матрицу метрики.

Для гладкой кривой $\gamma:[a,b]\to M$ её скорость принадлежит $T_{\gamma(t)}M$, а длина определяется формулой
\begin{equation}
\label{eq:riemannian-curve-length}
L_g(\gamma)
=
\int_a^b
\sqrt{g_{\gamma(t)}\bigl(\dot\gamma(t),\dot\gamma(t)\bigr)}\,\dd t.
\end{equation}
В координатах $x(t)=(x^1(t),\ldots,x^n(t))$ подынтегральное выражение имеет вид
\[
\sqrt{\dot x(t)^{\trans}g(x(t))\dot x(t)}.
\]
Формула не зависит от параметризации кривой, поскольку множитель изменения скорости сокращается с заменой переменной в интеграле.

Метрика задаёт и естественную форму объёма. В положительно ориентированной карте
\begin{equation}
\label{eq:riemannian-volume-form}
\dd V_g
=
\sqrt{\det g(x)}\,
\dd x^1\wedge\cdots\wedge\dd x^n.
\end{equation}
Определитель измеряет квадрат коэффициента изменения $n$-мерного объёма, поэтому в формуле появляется корень. При $n=2$ получаем элемент площади
\[
\dd A_g=\sqrt{\det g}\,\dd x^1\dd x^2,
\]
что совпадает с формулой $\sqrt{EG-F^2}\,\dd u\dd v$ для параметризованной поверхности. Ориентация из раздела~\ref{sec:differential-forms-manifolds} определяет знак формы, а метрика определяет её плотность.

\subsection{Риманов градиент}

Пусть $f:M\to\R$ --- гладкая функция. Её дифференциал
\[
\dd f_p:T_pM\to\R
\]
является ковектором и не зависит от выбора метрики. Чтобы превратить этот ковектор в направление движения, требуется скалярное произведение.

\begin{definition}
Римановым градиентом функции $f$ в точке $p$ называется единственный вектор $\grad_g f(p)\in T_pM$, удовлетворяющий равенству
\begin{equation}
\label{eq:riemannian-gradient-definition}
g_p\bigl(\grad_g f(p),v\bigr)
=
\dd f_p(v)
\qquad
\text{для всех }v\in T_pM.
\end{equation}
\end{definition}

В координатах положим
\[
\partial f
=
\begin{pmatrix}
\pdv{f}{x^1}&\cdots&\pdv{f}{x^n}
\end{pmatrix}^{\trans}.
\]
Если столбец коэффициентов градиента обозначить через $q$, то из определения следует
\[
q^{\trans}gv=(\partial f)^{\trans}v
\qquad\text{для всех }v.
\]
Следовательно,
\begin{equation}
\label{eq:riemannian-gradient-coordinates}
\boxed{\grad_g f=g^{-1}\partial f}.
\end{equation}
При $g=I$ получается обычный евклидов градиент. Метрика меняет направление наискорейшего роста: один и тот же дифференциал превращается в разные векторы при разных правилах измерения длины.

Для единичного касательного вектора $v$ производная по направлению равна
\[
\dd f_p(v)=g_p(\grad_g f,v)
\leq\norm{\grad_g f}_{g,p}.
\]
Равенство достигается при
\[
v=\frac{\grad_g f}{\norm{\grad_g f}_{g,p}}.
\]
Поэтому $-\grad_g f$ является направлением наиболее быстрого локального убывания функции среди касательных направлений единичной длины.

\subsection{Касательная проекция и ограничения}

Для вложенного многообразия $M\subseteq\R^m$ с индуцированной метрикой риманов градиент получается ортогональной проекцией обычного градиента на касательное пространство:
\begin{equation}
\label{eq:embedded-riemannian-gradient-projection}
\grad_M f(p)
=
P_{T_pM}\grad \widetilde f(p),
\end{equation}
где $\widetilde f$ --- любое гладкое продолжение $f$ в окрестность $M$. Нормальная компонента не влияет на производные вдоль допустимых направлений и потому удаляется проекцией.

Если многообразие задано системой ограничений
\[
M=\left\{x\in\R^m\middle|F(x)=0\right\},
\qquad
F:\R^m\to\R^k,
\]
и $DF(p)$ имеет ранг $k$, то
\[
T_pM=\ker DF(p).
\]
Ортогональная проекция на это ядро имеет вид
\begin{equation}
\label{eq:constraint-tangent-projector}
P_{T_pM}
=
I-DF(p)^{\trans}
\bigl(DF(p)DF(p)^{\trans}\bigr)^{-1}DF(p).
\end{equation}
Матрица в скобках имеет размер $k\times k$ и обратима благодаря условию полного ранга. Формула сначала находит нормальную компоненту в линейной оболочке градиентов ограничений, затем вычитает её.

В критической точке ограниченной задачи
\[
\grad_M f(p)=0.
\]
Это означает, что обычный градиент лежит в нормальном пространстве:
\begin{equation}
\label{eq:lagrange-geometric-form}
\grad\widetilde f(p)
=
DF(p)^{\trans}\lambda
\end{equation}
для некоторого $\lambda\in\R^k$. Так геометрическое условие нулевого касательного градиента совпадает с методом множителей Лагранжа из многомерного анализа.

\subsection{Шаг оптимизации и ретракция}

Пусть $p_k\in M$. Формальный градиентный шаг
\[
p_k-\eta_k\grad_M f(p_k)
\]
выполняется в окружающем линейном пространстве и обычно уже не принадлежит $M$. Поэтому касательное смещение возвращают на многообразие специальным отображением.

\begin{definition}
Ретракцией называется гладкое отображение
\[
R_p:T_pM\to M,
\]
которое в нуле удовлетворяет условиям
\[
R_p(0)=p,
\qquad
DR_p(0)=\operatorname{Id}_{T_pM}.
\]
\end{definition}

Первое условие сохраняет исходную точку, второе означает, что для малых $\xi\in T_pM$ кривая $t\mapsto R_p(t\xi)$ имеет начальную скорость $\xi$. Первый шаг риманова градиентного метода записывается как
\begin{equation}
\label{eq:riemannian-gradient-step}
p_{k+1}
=
R_{p_k}\bigl(-\eta_k\grad_M f(p_k)\bigr).
\end{equation}
Ретракция не обязана возвращать ближайшую точку и не обязана совпадать с движением по геодезической. Для алгоритма достаточно правильного поведения в первом порядке.

На сфере радиуса $R$
\[
S_R^{m-1}=\left\{x\in\R^m\middle|x^{\trans}x=R^2\right\}
\]
касательное пространство равно
\[
T_pS_R^{m-1}=\left\{v\in\R^m\middle|p^{\trans}v=0\right\}.
\]
Проектор и удобная ретракция имеют вид
\begin{equation}
\label{eq:sphere-projector-retraction}
P_p=I-\frac{pp^{\trans}}{R^2},
\qquad
R_p(\xi)=R\frac{p+\xi}{\norm{p+\xi}}.
\end{equation}
Нормировка возвращает точку на сферу, а при $\xi=0$ даёт $p$. Поскольку $p^{\trans}\xi=0$, производная нормировки в нуле оставляет касательное направление неизменным.

\subsection{Квадратичная функция на сфере}

Рассмотрим симметричную матрицу $A\in\R^{m\times m}$ и задачу
\begin{equation}
\label{eq:rayleigh-sphere-problem}
\min_{x\in S^{m-1}}
f(x),
\qquad
f(x)=\frac12x^{\trans}Ax.
\end{equation}
Евклидов градиент равен $\grad f(x)=Ax$. Для единичной сферы проектор из \eqref{eq:sphere-projector-retraction} даёт
\begin{equation}
\label{eq:rayleigh-riemannian-gradient}
\grad_S f(x)
=
Ax-(x^{\trans}Ax)x.
\end{equation}
В критической точке
\[
Ax=(x^{\trans}Ax)x,
\]
то есть $x$ является собственным вектором $A$, а число $x^{\trans}Ax$ --- соответствующим собственным значением. Минимум достигается на собственном направлении с наименьшим собственным значением. Множитель Лагранжа и геометрическая нормальная компонента здесь являются одним и тем же числом.

Пусть
\[
A=\diagop(1,2,4),
\qquad
x_0=\frac{1}{\sqrt3}(1,1,1)^{\trans},
\qquad
\eta=\frac13.
\]
Сначала вычисляется
\[
x_0^{\trans}Ax_0=\frac73,
\]
а затем касательный градиент
\begin{align}
\grad_S f(x_0)
&=Ax_0-\frac73x_0\notag\\
&=\frac{1}{3\sqrt3}(-4,-1,5)^{\trans}.
\label{eq:sphere-gradient-numerical}
\end{align}
Он ортогонален $x_0$, поскольку сумма произведений соответствующих координат равна нулю. Ненормированный шаг равен
\[
y_1
=x_0-\frac13\grad_S f(x_0)
=\frac{1}{9\sqrt3}(13,10,4)^{\trans}.
\]
После ретракции получаем
\begin{equation}
\label{eq:sphere-retracted-step-numerical}
x_1
=\frac{y_1}{\norm{y_1}}
=\frac{1}{\sqrt{285}}(13,10,4)^{\trans}.
\end{equation}
Значение функции уменьшилось:
\[
f(x_0)=\frac76\approx1.167,
\qquad
f(x_1)=\frac{433}{570}\approx0.760.
\]
Одна итерация состоит из трёх операций: вычисления обычного градиента, удаления нормальной компоненты и нормировки результата. Повторение шага при подходящем выборе длины шага приближает вектор к направлению $(1,0,0)^{\trans}$, соответствующему наименьшему собственному значению. Этот пример показывает, как линейная алгебра, множители Лагранжа и градиентный метод объединяются в одной геометрической схеме.

\section{Связность, параллельный перенос и геодезические}
\label{sec:connections-geodesics}

Вектор из $T_pM$ нельзя напрямую вычесть из вектора из $T_qM$: эти векторы принадлежат разным линейным пространствам. По той же причине обычная производная векторного поля вдоль кривой не имеет внутреннего смысла без дополнительного правила сравнения соседних касательных пространств. Такое правило задаёт связность. Она позволяет дифференцировать векторные поля, переносить направления вдоль кривых и формулировать условие движения ``без поворота'' на искривлённом пространстве.

\subsection{Аффинная связность}

Обозначим через $\mathfrak X(M)$ пространство гладких векторных полей на многообразии $M$.

\begin{definition}
Аффинной связностью на $M$ называется отображение
\[
\nabla:\mathfrak X(M)\times\mathfrak X(M)\to\mathfrak X(M),
\qquad
(X,Y)\mapsto\nabla_XY,
\]
линейное над $\R$ по обоим аргументам и удовлетворяющее правилам
\begin{equation}
\label{eq:affine-connection-rules}
\nabla_{fX}Y=f\nabla_XY,
\qquad
\nabla_X(fY)=X(f)Y+f\nabla_XY
\end{equation}
для любой гладкой функции $f:M\to\R$.
\end{definition}

Величина $\nabla_XY$ называется ковариантной производной поля $Y$ в направлении $X$. Первое равенство показывает, что в точке $p$ результат зависит только от вектора $X(p)$. Второе является правилом Лейбница: изменение произведения $fY$ складывается из изменения коэффициента $f$ и изменения самого поля $Y$.

Пусть $(x^1,\ldots,x^n)$ --- локальные координаты, а
\[
\partial_i:=\pdv{}{x^i}
\]
--- координатные поля. Связность полностью определяется равенствами
\begin{equation}
\label{eq:christoffel-definition}
\nabla_{\partial_j}\partial_k
=
\sum_{i=1}^n\Gamma^i_{jk}\partial_i.
\end{equation}
Функции $\Gamma^i_{jk}$ называются символами Кристоффеля. Если
\[
X=\sum_{j=1}^nX^j\partial_j,
\qquad
Y=\sum_{k=1}^nY^k\partial_k,
\]
то правила~\eqref{eq:affine-connection-rules} дают
\begin{equation}
\label{eq:covariant-derivative-coordinates}
(\nabla_XY)^i
=
\sum_{j=1}^nX^j\pdv{Y^i}{x^j}
+
\sum_{j,k=1}^n\Gamma^i_{jk}X^jY^k.
\end{equation}
Первое слагаемое совпадает с обычной производной координат компонент. Второе компенсирует изменение самого координатного базиса. Поэтому символы Кристоффеля зависят от карты и не являются компонентами тензора. Даже в евклидовой плоскости они равны нулю в декартовых координатах, но становятся ненулевыми в полярных.

\subsection{Связность Леви--Чивиты}

На римановом многообразии связность должна быть согласована с метрикой. Для полей $X,Y,Z$ метрическая совместимость записывается как
\begin{equation}
\label{eq:metric-compatibility}
Z\bigl(g(X,Y)\bigr)
=
g(\nabla_ZX,Y)+g(X,\nabla_ZY).
\end{equation}
Это обычное правило дифференцирования скалярного произведения, перенесённое на многообразие.

Кручением связности называется поле
\begin{equation}
\label{eq:connection-torsion-definition}
T(X,Y)=\nabla_XY-\nabla_YX-[X,Y].
\end{equation}
Условие $T=0$ означает, что ковариантные производные в двух направлениях отличаются ровно на скобку Ли. В координатном базисе $[\partial_j,\partial_k]=0$, поэтому отсутствие кручения равносильно симметрии
\[
\Gamma^i_{jk}=\Gamma^i_{kj}.
\]

\begin{theorem}[основная теорема римановой геометрии]
Для каждой римановой метрики $g$ существует единственная аффинная связность, которая совместима с $g$ и не имеет кручения. Она называется связностью Леви--Чивиты.
\end{theorem}

Таким образом, для римановой геометрии связность не является новым независимым выбором: она однозначно восстанавливается по метрике. В координатах этот факт выражается формулой
\begin{equation}
\label{eq:christoffel-from-metric}
\boxed{
\Gamma^i_{jk}
=
\frac12
\sum_{\ell=1}^n
 g^{i\ell}
\left(
\pdv{g_{k\ell}}{x^j}
+
\pdv{g_{j\ell}}{x^k}
-
\pdv{g_{jk}}{x^\ell}
\right),
}
\end{equation}
где $(g^{i\ell})=g^{-1}$. Для вычисления сначала строят матрицу $g=(g_{ij})$, затем находят её производные и обратную матрицу. На выходе получается $n^3$ функций $\Gamma^i_{jk}$, причём из симметрии по нижним индексам достаточно хранить только случаи $j\leq k$.

\subsection{Ковариантная производная вдоль кривой}

Пусть $\gamma:I\to M$ --- гладкая кривая, а
\[
V(t)\in T_{\gamma(t)}M
\]
--- векторное поле вдоль неё. В координатах
\[
\gamma(t)=\bigl(x^1(t),\ldots,x^n(t)\bigr),
\qquad
V(t)=\sum_{i=1}^nV^i(t)\partial_i.
\]
Ковариантная производная вдоль кривой определяется формулой
\begin{equation}
\label{eq:covariant-derivative-along-curve}
\frac{DV}{\dd t}
:=
\nabla_{\dot\gamma}V
=
\sum_{i=1}^n
\left(
\dv{V^i}{t}
+
\sum_{j,k=1}^n
\Gamma^i_{jk}(x(t))\dot x^jV^k
\right)\partial_i.
\end{equation}
Здесь входом являются координаты точки $x(t)$, скорость $\dot x(t)$ и компоненты $V(t)$; выход принадлежит тому же пространству $T_{\gamma(t)}M$, что и $V(t)$.

Если $M\subseteq\R^m$ вложено и метрика индуцирована из окружающего пространства, та же производная имеет наглядный вид
\begin{equation}
\label{eq:embedded-covariant-derivative}
\frac{DV}{\dd t}
=
P_{T_{\gamma(t)}M}\frac{\dd V}{\dd t}.
\end{equation}
Обычная производная $\dd V/\dd t$ может иметь нормальную составляющую; связность Леви--Чивиты оставляет только касательную часть. Формула~\eqref{eq:embedded-covariant-derivative} зависит от вложения в записи, но результат определяется внутренней метрикой.

\subsection{Параллельный перенос}

\begin{definition}
Поле $V(t)$ называется параллельным вдоль $\gamma$, если
\begin{equation}
\label{eq:parallel-field-equation}
\frac{DV}{\dd t}=0.
\end{equation}
\end{definition}

В координатах это линейная система обыкновенных дифференциальных уравнений
\begin{equation}
\label{eq:parallel-transport-ode}
\dv{V^i}{t}
+
\sum_{j,k=1}^n
\Gamma^i_{jk}(x(t))\dot x^jV^k
=0,
\qquad i=1,\ldots,n.
\end{equation}
Начальный вектор $V(t_0)\in T_{\gamma(t_0)}M$ однозначно определяет решение и тем самым линейное отображение
\[
P_{t_0\to t}:T_{\gamma(t_0)}M\to T_{\gamma(t)}M,
\]
которое называется параллельным переносом вдоль $\gamma$. Из метрической совместимости следует
\[
\dv{}{t}g\bigl(V(t),W(t)\bigr)=0
\]
для любых параллельных полей $V,W$. Следовательно, перенос сохраняет длины и углы. Он даёт корректный способ сравнить направления в разных точках, но результат обычно зависит от выбранного пути. Эта зависимость станет одним из проявлений кривизны в следующем разделе.

\subsection{Геодезические}

В евклидовом пространстве прямая характеризуется постоянством вектора скорости. На многообразии обычная производная скорости заменяется ковариантной.

\begin{definition}
Кривая $\gamma$ называется геодезической, если её скорость параллельна вдоль самой кривой:
\begin{equation}
\label{eq:geodesic-invariant-equation}
\frac{D\dot\gamma}{\dd t}=0.
\end{equation}
\end{definition}

Подстановка $V=\dot\gamma$ в~\eqref{eq:covariant-derivative-along-curve} даёт геодезическое уравнение
\begin{equation}
\label{eq:geodesic-coordinate-equation}
\boxed{
\ddot x^i
+
\sum_{j,k=1}^n
\Gamma^i_{jk}(x)\dot x^j\dot x^k
=0,
\qquad i=1,\ldots,n.
}
\end{equation}
Это система $n$ дифференциальных уравнений второго порядка. Начальные данные
\[
x(0)=x_0,
\qquad
\dot x(0)=v_0
\]
задают локально единственную геодезическую. Для численного решения систему переписывают как $2n$ уравнений первого порядка для пары $(x,v)$:
\begin{equation}
\label{eq:geodesic-first-order-system}
\dot x^i=v^i,
\qquad
\dot v^i=-\sum_{j,k=1}^n\Gamma^i_{jk}(x)v^jv^k.
\end{equation}
Поскольку связность совместима с метрикой,
\[
\dv{}{t}g(\dot\gamma,\dot\gamma)
=2g\left(\frac{D\dot\gamma}{\dd t},\dot\gamma\right)=0.
\]
Значит, геодезическая имеет постоянную скорость. Параметр $t$ здесь аффинный: линейная замена $t\mapsto at+b$ сохраняет уравнение, а произвольная перепараметризация обычно нет.

\subsection{Энергия и кратчайшие пути}

Для кривой $\gamma:[a,b]\to M$ определим энергию
\begin{equation}
\label{eq:curve-energy}
\mathcal E(\gamma)
=
\frac12\int_a^b
 g_{\gamma(t)}(\dot\gamma(t),\dot\gamma(t))\,\dd t
=
\frac12\int_a^b
\sum_{i,j=1}^n g_{ij}(x(t))\dot x^i\dot x^j\,\dd t.
\end{equation}
Геодезические являются критическими точками этой энергии среди гладких вариаций с фиксированными концами. Уравнения Эйлера--Лагранжа для подынтегральной функции в~\eqref{eq:curve-energy} после умножения на $g^{-1}$ приводят именно к~\eqref{eq:geodesic-coordinate-equation}.

Длина и энергия связаны неравенством Коши--Буняковского
\begin{equation}
\label{eq:length-energy-inequality}
L_g(\gamma)^2
\leq
2(b-a)\mathcal E(\gamma),
\end{equation}
причём равенство достигается при постоянной скорости. Поэтому локально кратчайшая кривая, параметризованная с постоянной скоростью, является геодезической. Обратное утверждение верно лишь локально: после достаточно длинного движения геодезическая может перестать быть кратчайшей. Например, две точки сферы могут соединяться несколькими дугами больших кругов, но минимальна только более короткая дуга.

Теорема Хопфа--Ринова связывает локальное уравнение с глобальной геометрией. Для связного риманова многообразия метрическая полнота равносильна возможности продолжать каждую геодезическую на все значения параметра. В этом случае любые две точки соединяются минимизирующей геодезической. Компактные римановы многообразия полны, поэтому на сфере такая геодезическая всегда существует.

\subsection{Практический расчёт: геодезические на цилиндре}

Рассмотрим цилиндр радиуса $R$ с параметризацией
\begin{equation}
\label{eq:cylinder-parameterization-v07}
X(\theta,z)
=
\bigl(R\cos\theta,R\sin\theta,z\bigr),
\qquad
(\theta,z)\in\R^2.
\end{equation}
Координатные касательные векторы равны
\[
X_\theta=(-R\sin\theta,R\cos\theta,0),
\qquad
X_z=(0,0,1),
\]
поэтому индуцированная метрика имеет постоянную матрицу
\begin{equation}
\label{eq:cylinder-metric-v07}
g=
\begin{pmatrix}
R^2&0\\
0&1
\end{pmatrix},
\qquad
g^{-1}=
\begin{pmatrix}
R^{-2}&0\\
0&1
\end{pmatrix}.
\end{equation}
Все производные $\partial g_{ij}/\partial x^k$ равны нулю. Из формулы~\eqref{eq:christoffel-from-metric} следует
\[
\Gamma^i_{jk}=0.
\]
Геодезические уравнения становятся
\[
\ddot\theta=0,
\qquad
\ddot z=0,
\]
и имеют решения
\begin{equation}
\label{eq:cylinder-geodesic-solution}
\theta(t)=\theta_0+at,
\qquad
z(t)=z_0+bt.
\end{equation}
В пространстве такая кривая имеет вид
\[
\gamma(t)
=
\bigl(R\cos(\theta_0+at),
R\sin(\theta_0+at),
z_0+bt\bigr).
\]
При $a=0$ получается образующая цилиндра, при $b=0$ --- окружность, а при $ab\neq0$ --- винтовая линия. Её скорость постоянна:
\[
\norm{\dot\gamma}^2=R^2a^2+b^2.
\]

Проверка не требует повторного дифференцирования в $\R^3$. Разрежем цилиндр вдоль образующей и введём координаты развёртки
\[
u=R\theta,
\qquad
v=z.
\]
Тогда
\[
\dd s^2=R^2\dd\theta^2+\dd z^2=\dd u^2+\dd v^2,
\]
а решение~\eqref{eq:cylinder-geodesic-solution} переходит в прямую
\[
(u(t),v(t))=(R\theta_0,z_0)+t(Ra,b).
\]
На участке, не пересекающем разрез и не меняющем число оборотов, прямая минимизирует длину, поэтому соответствующая винтовая линия действительно является локально кратчайшей на цилиндре.

Тот же расчёт объясняет параллельный перенос. Поскольку $\Gamma^i_{jk}=0$, уравнение~\eqref{eq:parallel-transport-ode} даёт
\[
\dot V^\theta=0,
\qquad
\dot V^z=0.
\]
Компоненты в базисе $(X_\theta,X_z)$ сохраняются, хотя сам вектор в окружающем $\R^3$ поворачивается вместе с цилиндром. Внутренняя геометрия считает его направление неизменным.

\subsection{Сфера: пример ненулевых символов}

Для сферы радиуса $R$ возьмём координаты
\[
X(\varphi,\lambda)
=
\bigl(R\sin\varphi\cos\lambda,
R\sin\varphi\sin\lambda,
R\cos\varphi\bigr),
\]
где $\varphi\in(0,\pi)$ --- полярный угол, а $\lambda$ --- долгота. Метрика равна
\begin{equation}
\label{eq:sphere-metric-v07}
g=
\begin{pmatrix}
R^2&0\\
0&R^2\sin^2\varphi
\end{pmatrix}.
\end{equation}
Здесь коэффициент $g_{\lambda\lambda}$ зависит от $\varphi$, поэтому часть символов Кристоффеля ненулевая. Подстановка в~\eqref{eq:christoffel-from-metric} даёт
\begin{equation}
\label{eq:sphere-christoffel-v07}
\Gamma^\varphi_{\lambda\lambda}
=-\sin\varphi\cos\varphi,
\qquad
\Gamma^\lambda_{\varphi\lambda}
=
\Gamma^\lambda_{\lambda\varphi}
=\cot\varphi.
\end{equation}
Следовательно, геодезическое уравнение принимает вид
\begin{equation}
\label{eq:sphere-geodesic-equations-v07}
\ddot\varphi-
\sin\varphi\cos\varphi\,\dot\lambda^2=0,
\qquad
\ddot\lambda+2\cot\varphi\,\dot\varphi\dot\lambda=0.
\end{equation}
Для экватора $\varphi(t)=\pi/2$ и $\lambda(t)=at+b$ обе левые части равны нулю. Для меридиана $\lambda=\mathrm{const}$ остаётся $\ddot\varphi=0$. Эти две проверки подтверждают, что большие круги являются геодезическими. В общем случае вычислительный конвейер остаётся тем же: параметризация задаёт $g$, производные метрики задают $\Gamma^i_{jk}$, а затем система~\eqref{eq:geodesic-first-order-system} интегрируется как обычная задача Коши для ОДУ.

\section{Кривизна и переход к дискретной геометрии}
\label{sec:curvature-discrete-geometry}

Метрика определяет длины и углы внутри поверхности, а связность --- способ сравнивать касательные векторы в соседних точках. Кривизна отвечает на следующий вопрос: насколько геометрия поверхности отличается от геометрии плоскости и как это отличие проявляется при движении, переносе направлений и изменении формы. Для гладкой поверхности кривизна выражается через производные нормали и связности. Для треугольной сетки производных в обычном смысле нет, поэтому те же величины заменяются конечными суммами по рёбрам, углам и граням.

\subsection{Вторая фундаментальная форма и оператор формы}

Пусть ориентированная поверхность задана регулярной параметризацией
\[
X:U\subseteq\R^2\to\R^3,
\qquad
(u,v)\mapsto X(u,v),
\]
а единичная нормаль выбрана по формуле
\[
n=\frac{X_u\times X_v}{\norm{X_u\times X_v}}.
\]
Первая фундаментальная форма из раздела~\ref{sec:parametric-surfaces} измеряет длины касательных векторов. Чтобы описать изгиб поверхности в окружающем пространстве, нужно проследить, как меняется нормаль при перемещении по поверхности.

Из равенства \(\inner{n}{n}=1\) следует
\[
0=D_w\inner{n}{n}=2\inner{D_wn}{n}
\]
для любого касательного направления \(w\in T_pS\). Поэтому производная \(D_wn\) снова лежит в касательной плоскости. Это позволяет определить линейный оператор на одном и том же пространстве.

\begin{definition}
Оператором формы поверхности в точке \(p\) называется отображение
\begin{equation}
\label{eq:shape-operator}
S_p:T_pS\to T_pS,
\qquad
S_p(w)=-D_wn.
\end{equation}
Второй фундаментальной формой называется симметричная билинейная форма
\begin{equation}
\label{eq:second-fundamental-form}
\mathrm{II}_p(w,z)=\inner{S_p(w)}{z}.
\end{equation}
\end{definition}

Знак в~\eqref{eq:shape-operator} является соглашением. При выбранном знаке замена ориентации \(n\mapsto-n\) меняет знаки \(S\), \(\mathrm{II}\) и средней кривизны, но не меняет гауссову кривизну.

В координатном базисе \((X_u,X_v)\) матрица первой фундаментальной формы равна
\begin{equation}
\label{eq:v08-first-form-matrix}
G=
\begin{pmatrix}
E&F\\
F&G_0
\end{pmatrix}
=
\begin{pmatrix}
\inner{X_u}{X_u}&\inner{X_u}{X_v}\\
\inner{X_u}{X_v}&\inner{X_v}{X_v}
\end{pmatrix}.
\end{equation}
Здесь нижний индекс у \(G_0\) нужен только для отличия коэффициента классической записи от всей матрицы \(G\). Дифференцируя равенства \(\inner{n}{X_u}=\inner{n}{X_v}=0\), получаем
\[
-\inner{n_u}{X_u}=\inner{n}{X_{uu}},
\qquad
-\inner{n_u}{X_v}=\inner{n}{X_{uv}},
\qquad
-\inner{n_v}{X_v}=\inner{n}{X_{vv}}.
\]
Поэтому матрица второй фундаментальной формы имеет вид
\begin{equation}
\label{eq:second-form-matrix}
B=
\begin{pmatrix}
e&f\\
f&g
\end{pmatrix},
\qquad
e=\inner{X_{uu}}{n},
\quad
f=\inner{X_{uv}}{n},
\quad
g=\inner{X_{vv}}{n}.
\end{equation}

Если координатный вектор \(a=(a^1,a^2)^{\trans}\) представляет касательный вектор \(w=a^1X_u+a^2X_v\), то
\[
\mathrm{I}(w,w)=a^{\trans}Ga,
\qquad
\mathrm{II}(w,w)=a^{\trans}Ba.
\]
Матрица оператора формы определяется равенством
\begin{equation}
\label{eq:shape-matrix-from-forms}
\boxed{[S]=G^{-1}B.}
\end{equation}
Матрица \(G^{-1}B\) не обязана быть симметричной относительно обычного скалярного произведения в \(\R^2\), но оператор \(S\) самосопряжён относительно метрики \(G\):
\[
a^{\trans}GSb=b^{\trans}GSa.
\]
Именно поэтому его собственные значения вещественны, а собственные направления можно выбрать взаимно ортогональными на поверхности.

\subsection{Нормальная, главные, средняя и гауссова кривизны}

Пусть \(w\neq0\) --- касательное направление. Плоскость, натянутая на \(w\) и нормаль \(n\), пересекает поверхность по кривой. Её кривизна в нормальном направлении называется нормальной кривизной поверхности вдоль \(w\). Через фундаментальные формы она вычисляется как отношение
\begin{equation}
\label{eq:normal-curvature-surface}
\kappa_n(w)
=
\frac{\mathrm{II}(w,w)}{\mathrm{I}(w,w)}
=
\frac{a^{\trans}Ba}{a^{\trans}Ga}.
\end{equation}
Масштаб вектора \(a\) сокращается, поэтому значение зависит только от направления.

\begin{definition}
Главными направлениями называются собственные направления оператора формы,
\[
S(e_i)=\kappa_i e_i,
\qquad i=1,2.
\]
Собственные значения \(\kappa_1,\kappa_2\) называются главными кривизнами. Средняя и гауссова кривизны определяются формулами
\begin{equation}
\label{eq:mean-gaussian-curvature}
H=\frac{\kappa_1+\kappa_2}{2}
=\frac12\operatorname{tr}S,
\qquad
K=\kappa_1\kappa_2
=\det S
=\frac{\det B}{\det G}.
\end{equation}
\end{definition}

Главные кривизны имеют размерность \(1/\text{длина}\). Средняя кривизна имеет ту же размерность и чувствительна к ориентации нормали. Гауссова кривизна имеет размерность \(1/\text{длина}^2\) и от ориентации не зависит.

Если \(K>0\), главные кривизны имеют одинаковый знак: поверхность локально похожа на купол или чашу. Если \(K<0\), знаки противоположны и поверхность имеет седловой характер. При \(K=0\) хотя бы одно главное направление плоское. У цилиндра одно главное значение равно нулю, поэтому цилиндр можно развернуть на плоскость без локального растяжения. У сферы обе главные кривизны по модулю равны \(1/R\); все направления главные, а \(K=1/R^2\).

Средняя кривизна описывает первую вариацию площади. Если поверхность сдвинуть с нормальной скоростью \(v_n\), то изменение площади в первом порядке содержит интеграл от \(H v_n\). Поэтому условие \(H=0\) возникает у минимальных поверхностей, а движение в направлении вектора средней кривизны сглаживает мелкие неровности.

\subsection{Внутренняя кривизна и тензор Римана}

Вторая фундаментальная форма использует нормаль и тем самым внешнее вложение поверхности в \(\R^3\). Гауссова кривизна, несмотря на формулу \(K=\det B/\det G\), определяется одной только метрикой. Это содержание теоремы Гаусса: наблюдатель, измеряющий лишь длины, углы и геодезические на поверхности, может восстановить \(K\), не обращаясь к окружающему пространству.

В произвольной размерности кривизна связности Леви--Чивиты кодируется тензором
\begin{equation}
\label{eq:riemann-curvature-tensor}
R(X,Y)Z
=
\nabla_X\nabla_YZ
-\nabla_Y\nabla_XZ
-\nabla_{[X,Y]}Z.
\end{equation}
Он измеряет некоммутативность ковариантного дифференцирования. Если малый вектор параллельно перенести по замкнутому контуру, то расхождение между начальным и конечным направлениями в первом приближении определяется площадью контура и значением \(R\). Для двумерной поверхности и ортонормированного базиса \((e_1,e_2)\) вся внутренняя кривизна сводится к одному числу
\begin{equation}
\label{eq:gaussian-from-riemann}
K=\inner{R(e_1,e_2)e_2}{e_1},
\end{equation}
с точностью до соглашения о знаке тензора \(R\). Полную координатную теорию компонент \(R^i{}_{jkl}\) здесь не развиваем: для дальнейших вычислений важнее понять, что кривизна проявляется как зависимость параллельного переноса от пути.

\subsection{Треугольная сетка как модель поверхности}

Компьютерная геометрия обычно хранит поверхность не формулой \(X(u,v)\), а конечным набором вершин
\[
p_1,\ldots,p_N\in\R^3,
\]
рёбер и ориентированных треугольных граней. Комбинаторная часть сообщает, какие вершины соединены; координаты задают геометрию. Окрестность вершины \(p_i\), состоящая из всех инцидентных ей треугольников, называется её кольцом первого порядка.

Внутри отдельной грани геометрия плоская: нормаль постоянна, а обычные производные второго порядка равны нулю. Кривизна сосредоточивается на стыках граней и в вершинах. Поэтому дискретный аналог гладкой величины чаще представляет не точечное значение, а её интеграл по небольшой окрестности. Для перехода к плотности этот интеграл делят на площадь соответствующей ячейки.

Пусть \(A_i\) --- площадь окрестности вершины. Простейший выбор --- одна треть суммы площадей всех прилегающих треугольников:
\begin{equation}
\label{eq:barycentric-vertex-area}
A_i=\frac13\sum_{f\ni i}\operatorname{Area}(f).
\end{equation}
Для более точных схем используют смешанные или вороновы ячейки, но принцип остаётся тем же: сначала вычисляется интегральная геометрическая величина, затем она нормируется на локальную площадь.

\subsection{Дефект угла и дискретная гауссова кривизна}

Пусть \(\theta_{if}\) --- внутренний угол грани \(f\) при внутренней вершине \(p_i\). В плоской триангуляции сумма углов вокруг вершины равна \(2\pi\). Отклонение от этого значения задаёт дефект угла
\begin{equation}
\label{eq:angle-defect}
\boxed{
\delta_i
=2\pi-\sum_{f\ni i}\theta_{if}.
}
\end{equation}
Величина \(\delta_i\) является дискретным аналогом интеграла гауссовой кривизны по окрестности вершины. Точечную оценку получают делением на площадь:
\begin{equation}
\label{eq:discrete-gaussian-curvature}
K_i\approx\frac{\delta_i}{A_i}.
\end{equation}
Положительный дефект соответствует куполообразной вершине, отрицательный --- седловой. Для вершины на границе вместо \(2\pi\) используется \(\pi\), поскольку её окрестность является половиной диска.

Для замкнутой ориентируемой треугольной поверхности выполняется точное дискретное равенство Гаусса--Бонне
\begin{equation}
\label{eq:discrete-gauss-bonnet}
\boxed{
\sum_{i=1}^N\delta_i
=2\pi\chi,
\qquad
\chi=N-E+F,
}
\end{equation}
где \(E\) и \(F\) --- числа рёбер и граней. Локальные дефекты меняются при деформации сетки, но их сумма определяется только топологией.

\begin{example}[регулярный тетраэдр]
У правильного тетраэдра \(N=4\), \(E=6\), \(F=4\), поэтому
\[
\chi=4-6+4=2.
\]
В каждой вершине сходятся три равносторонних треугольника, и каждый угол равен \(\pi/3\). Следовательно,
\[
\delta_i
=2\pi-3\frac{\pi}{3}
=\pi.
\]
Суммарный дефект равен \(4\pi\), что точно совпадает с \(2\pi\chi=4\pi\). При этом значение \(K_i=\pi/A_i\) зависит от масштаба тетраэдра, а интегральная кривизна \(\delta_i\) --- нет.
\end{example}

\subsection{Дискретный лапласиан и средняя кривизна}

Пусть скалярная функция задана значениями \(u_i\) в вершинах сетки и линейно интерполируется внутри каждой грани. Для внутреннего ребра \((i,j)\) обозначим через \(\alpha_{ij}\) и \(\beta_{ij}\) углы, лежащие напротив этого ребра в двух соседних треугольниках. Котангенсный вес ребра равен
\begin{equation}
\label{eq:cotangent-weight}
w_{ij}
=\frac12\bigl(\cot\alpha_{ij}+\cot\beta_{ij}\bigr).
\end{equation}
На граничном ребре остаётся только один угол.

Интегральный дискретный лапласиан в вершине определяется формулой
\begin{equation}
\label{eq:integrated-cotan-laplacian}
(Lu)_i
=\sum_{j\sim i}w_{ij}(u_j-u_i),
\end{equation}
где сумма берётся по соседним вершинам. После деления на площадь получается приближение оператора Лапласа--Бельтрами:
\begin{equation}
\label{eq:pointwise-cotan-laplacian}
(\Delta u)_i
\approx
\frac{1}{A_i}
\sum_{j\sim i}w_{ij}(u_j-u_i).
\end{equation}
Формула использует только локальные длины и углы. Матрица \(L\in\R^{N\times N}\) разрежена: ненулевые внедиагональные элементы возникают только для соседних вершин. Симметрия \(w_{ij}=w_{ji}\) делает \(L\) симметричной, а сумма элементов каждой строки равна нулю. Поэтому константы лежат в ядре лапласиана.

Применим оператор покомпонентно к вектору координат
\[
p=(p_1,\ldots,p_N)^{\trans},
\qquad p_i\in\R^3.
\]
Гладкое тождество \(\Delta X=2Hn\) переходит в приближение
\begin{equation}
\label{eq:discrete-mean-curvature-normal}
\frac{1}{A_i}(Lp)_i
\approx 2H_i n_i.
\end{equation}
Таким образом, одна и та же котангенсная сумма одновременно даёт направление нормали, вектор средней кривизны и градиент площади с точностью до выбранных знаков и нормировок. Это важный принцип дискретной дифференциальной геометрии: дискретизируется не отдельная формула, а согласованная система связей между объектами.

Котангенсные веса могут быть отрицательными на сетке с тупыми углами. Это не является алгебраической ошибкой, но способно ухудшить монотонность и устойчивость некоторых алгоритмов. На практике контролируют качество треугольников, используют делонеевскую перестройку рёбер или выбирают другую дискретизацию, если знакоположительность весов критична.

\subsection{Уравнение Пуассона на сетке}

Для уравнения
\[
\Delta u=f
\]
введём диагональную матрицу масс
\begin{equation}
\label{eq:mesh-mass-matrix}
M=\operatorname{diag}(A_1,\ldots,A_N).
\end{equation}
С учётом интегральной формы~\eqref{eq:integrated-cotan-laplacian} дискретная задача принимает вид
\begin{equation}
\label{eq:mesh-poisson-system}
Lu=Mf.
\end{equation}
На замкнутой связной поверхности решение существует лишь при нулевом среднем правой части и определяется с точностью до добавления константы. Для устранения неоднозначности фиксируют значение в одной вершине либо добавляют условие
\[
\sum_{i=1}^NA_i u_i=0.
\]
Если поверхность имеет край и значения \(u\) заданы на граничных вершинах, неизвестными остаются только внутренние значения. После разбиения вершин на внутренние \(I\) и граничные \(B\) система записывается как
\begin{equation}
\label{eq:dirichlet-block-system}
L_{II}u_I
=(Mf)_I-L_{IB}u_B.
\end{equation}
Это обычная разреженная линейная система. Геометрическая задача тем самым сводится к матрицам из главы линейной алгебры.

Уравнение Пуассона используется не только для скалярных физических полей. При параметризации поверхности две функции \(u\) и \(v\) задают координаты развёртки, а их внутренние значения находят из двух лапласовских или пуассоновских задач с закреплённой границей. При проектировании касательных векторных полей разложение Ходжа также приводит к скалярным уравнениям Пуассона для потенциалов. Детали этих алгоритмов зависят от граничных условий и топологии, но вычислительное ядро остаётся тем же.

\subsection{Сглаживание и поток средней кривизны}

Пусть положение поверхности зависит от времени: \(p=p(t)\). Уравнение
\begin{equation}
\label{eq:mean-curvature-flow-smooth}
\pdv{p}{t}=\Delta p
\end{equation}
двигает точки в направлении вектора средней кривизны. Выпуклые выступы и мелкие колебания уменьшаются, поэтому сетка сглаживается. Явная схема
\[
p^{k+1}=p^k+hM^{-1}Lp^k
\]
проста, но требует малого шага \(h\). Неявная схема использует лапласиан в новом состоянии:
\begin{equation}
\label{eq:implicit-mean-curvature-flow}
\boxed{
(M-hL)p^{k+1}=Mp^k.
}
\end{equation}
Матрица слева разрежена и одинакова для трёх координат, поэтому её факторизацию можно использовать трижды. Неявный шаг устойчивее и допускает существенно большие значения \(h\).

Сглаживание не является нейтральной операцией: поток средней кривизны обычно уменьшает объём, стирает острые рёбра и может сдвигать границу. Если требуется сохранить характерные линии, объём или закреплённые точки, эти условия нужно добавить в линейную систему или в последующую проекцию. Геометрический оператор задаёт направление изменения, а инженерные ограничения определяют допустимую часть этого изменения.

\subsection{Вычислительный конвейер}

Связь гладкой и дискретной геометрии можно записать как единый рабочий процесс:
\begin{equation}
\label{eq:geometry-processing-pipeline}
\boxed{
\begin{aligned}
&\text{параметризация или треугольная сетка}
\\[-0.2em]
&\quad\longrightarrow
\text{касательные данные, углы и локальные площади}
\\[-0.2em]
&\quad\longrightarrow
\text{метрика, кривизна, матрицы }L\text{ и }M
\\[-0.2em]
&\quad\longrightarrow
\text{линейная система или задача оптимизации}
\\[-0.2em]
&\quad\longrightarrow
\text{обновлённая форма, параметризация или поле}.
\end{aligned}
}
\end{equation}

На гладкой поверхности производные \(X_u,X_v,X_{uu},X_{uv},X_{vv}\) дают матрицы \(G\) и \(B\), затем оператор \(S=G^{-1}B\) и кривизны. На сетке их роль выполняют рёбра, углы и площади; из них собираются дефекты углов, котангенсный лапласиан и матрица масс. После этого задача становится знакомой: решить разреженную систему, минимизировать квадратичную энергию или выполнить шаг по времени. Именно здесь дифференциальная геометрия соединяется с линейной алгеброй, многомерным анализом, оптимизацией, обыкновенными и частными дифференциальными уравнениями.
